\documentclass[UTF8,zihao=5,a4paper]{ctexart}
\usepackage[margin=2.5cm]{geometry}
\usepackage{setspace}
\usepackage{booktabs}
\usepackage{array}
\usepackage{longtable}
\usepackage{tabularx}
\usepackage{amsmath,amssymb}
\usepackage{siunitx}
\usepackage{hyperref}
\usepackage{xcolor}
\usepackage{enumitem}
\usepackage{caption}
\usepackage{fancyhdr}
\usepackage{titlesec}
\usepackage{graphicx}
\usepackage{float}
\usepackage{tikz}
\usetikzlibrary{arrows.meta,positioning,shapes.geometric,fit,calc,backgrounds}

\hypersetup{colorlinks=true, linkcolor=black, citecolor=black, urlcolor=blue}
\sisetup{detect-weight=true,table-number-alignment=center}
\definecolor{tseBlue}{HTML}{DCEBFA}
\definecolor{tseBlueLine}{HTML}{3D6FA6}
\definecolor{tseGreen}{HTML}{E3F3E6}
\definecolor{tseGreenLine}{HTML}{3F8C5A}
\definecolor{tseAmber}{HTML}{FFF0D6}
\definecolor{tseAmberLine}{HTML}{B57A1B}
\definecolor{tseRed}{HTML}{FCE2E0}
\definecolor{tseRedLine}{HTML}{A94A45}
\definecolor{tseGray}{HTML}{F2F4F6}
\definecolor{tseGrayLine}{HTML}{667085}
\pagestyle{fancy}
\fancyhf{}
\cfoot{\thepage}
\renewcommand{\headrulewidth}{0pt}
\setstretch{1.35}
\setlength{\parindent}{2em}
\setlist[itemize]{leftmargin=2em,itemsep=0.2em,topsep=0.2em}
\setlist[enumerate]{leftmargin=2em,itemsep=0.2em,topsep=0.2em}
\captionsetup{font=small,labelsep=quad}
\titleformat{\section}{\zihao{4}\songti}{\thesection}{1em}{}
\titleformat{\subsection}{\zihao{5}\heiti}{\thesubsection}{1em}{}
\titleformat{\subsubsection}{\zihao{5}\heiti}{\thesubsubsection}{1em}{}
\makeatletter
\renewcommand{\@biblabel}[1]{#1.}
\makeatother

\newcommand{\sisdr}{SI-SDR}

\begin{document}

\thispagestyle{empty}
\vspace*{3.1cm}
\begin{center}
{\zihao{0}\kaishu 课程报告\par}
\end{center}
\vspace{4.2cm}
\begin{center}
\zihao{4}\songti
\renewcommand{\arraystretch}{1.45}
\begin{tabular}{rl}
姓名： & \underline{\makebox[8.0cm][c]{王志翔}}\\
学号： & \underline{\makebox[8.0cm][c]{22920242203422}}\\
年级： & \underline{\makebox[8.0cm][c]{2024级}}\\
专业方向： & \underline{\makebox[8.0cm][c]{人工智能专业}}\\
任课教师： & \underline{\makebox[8.0cm][c]{吴梅红}}\\
\end{tabular}
\end{center}
\vspace{2.0cm}
\begin{center}
{\zihao{-4}\songti 2026 年 6 月}\par
\end{center}
\newpage
\setcounter{page}{1}

\section{任务一：人工智能与人类智能的多维度对比分析}

\subsection{问题引入：为什么不能只用“强弱”评价智能}

人工智能与人类智能的比较，不能简单归结为“谁更聪明”。二者虽然都表现出解决问题、学习规律和作出决策的能力，但其实现机制、知识来源、泛化方式、目标系统和适用边界存在根本差异。当前人工智能，尤其是深度学习和大语言模型，在若干封闭任务中已经达到甚至超过普通人类水平，例如图像分类、语音识别、机器翻译、蛋白质结构预测、程序生成和标准化考试问答等。然而，人类智能并不只是题库上的高分表现，还包括具身感知、常识推理、社会理解、情绪调节、长期目标规划、自我意识和价值判断等复杂能力。

本节主要从认知架构、信息处理、学习泛化、因果推理、情感社会认知、具身性和评测边界几个方面比较二者差异。人工智能在\textbf{高维模式识别、大规模知识检索、形式化推理辅助和特定领域优化}上优势明显；人类智能在\textbf{开放世界理解、低样本学习、因果建模、价值判断、社会互动和具身适应}方面仍更稳定。

\subsection{研究问题与分析框架}

为了避免把人工智能和人类智能的差异简化为零散优缺点，本文采用“能力表现—实现机制—适用边界”的三层分析框架。第一层关注外显表现，即 AI 或人类在某项任务中的成绩，例如分类准确率、考试分数、推理正确率、任务完成时间和创造性产出数量。第二层关注实现机制，即这种表现依赖什么样的信息加工方式，是统计学习、符号推理、具身经验、社会互动，还是多种机制的结合。第三层关注适用边界，即这种能力在何种条件下有效，又会在哪些分布外、价值冲突或现实情境中失效。

这个框架可以避免把“表现相似”直接当成“机制相同”。例如，大语言模型可以写出看似有同理心的安慰语句，人类也可以表达同理心；但前者主要依赖语言模式生成，后者则与情绪体验、社会关系、责任意识和个人记忆相关。再如，AI 可以在某些数学竞赛题上取得高分，但这并不直接说明它具有人类数学家的问题意识、审美判断和长期研究动机。因此，本文讨论“超越”时只在具体任务和明确评价标准下使用该词，而不将单项任务高分等同于整体智能优越。

\begin{figure}[htbp]
\centering
\begin{tikzpicture}[
  node distance=1.0cm,
  box/.style={draw, rounded corners, align=center, minimum width=3.2cm, minimum height=0.85cm, font=\small},
  arrow/.style={-{Stealth[length=2.2mm]}, thick}
]
\node[box] (perf) {能力表现\\任务成绩、速度、准确率};
\node[box, right=of perf] (mechanism) {实现机制\\统计学习、符号推理、具身经验};
\node[box, right=of mechanism] (boundary) {适用边界\\分布外、伦理、责任};
\node[box, below=of mechanism, minimum width=4.0cm] (coop) {结论：人机协同\\AI 扩展计算能力，人类承担目标和价值};
\draw[arrow] (perf) -- (mechanism);
\draw[arrow] (mechanism) -- (boundary);
\draw[arrow] (boundary.south) |- (coop.east);
\draw[arrow] (perf.south) |- (coop.west);
\end{tikzpicture}
\caption{人工智能与人类智能比较的三层分析框架}
\label{fig:ai_human_framework}
\end{figure}

图\ref{fig:ai_human_framework}给出了本文的总体分析路径。先比较外显能力，再追问能力背后的实现机制，最后讨论这种能力能否迁移到开放环境和社会责任场景中。这样的顺序可以避免两种常见误区：一种是只看 benchmark 分数而忽略机制差异，另一种是只强调人类独特性而否认 AI 在具体任务中的真实进步。

\subsection{相关研究脉络综述}

人工智能与人类智能的比较经历了从行为标准、符号认知、联结主义到具身认知和大模型研究的多次转向。Turing 提出的“机器能否思考”问题把智能评价转化为可观察行为与交互表现\cite{turing1950computing}；Newell 和 Simon 的符号主义研究强调问题求解、规则和搜索过程\cite{newell1972human}；联结主义则把认知解释为大量简单单元连接权重变化的结果。进入深度学习时代后，人工智能在图像、语音、语言和科学计算中表现出强大模式学习能力，但具身认知研究提醒人们，真实智能还与身体行动、环境反馈和社会互动密切相关\cite{varela1991embodied,clark1997being}。

这些研究路线分别强调智能的不同侧面。符号主义有助于解释逻辑推理和显式规则，联结主义适合解释表示学习和模式识别，具身认知强调行动与环境，大模型研究展示了规模化数据和参数带来的能力跃迁。比较 AI 与人类智能时，不能只看模型规模，还要看知识来源、目标形成、不确定情境处理和行动后果承担。

\begin{longtable}{p{2.8cm}p{3.0cm}p{4.2cm}p{4.0cm}}
\toprule
研究传统 & 代表问题 & 对智能的理解 & 对本文比较的启发 \\
\midrule
\endfirsthead
\toprule
研究传统 & 代表问题 & 对智能的理解 & 对本文比较的启发 \\
\midrule
\endhead
行为评价传统 & 机器能否在交互中表现得像人 & 智能可通过外显行为和任务表现评价 & 说明 AI 可以在具体表现上接近或超过人 \\
符号主义传统 & 如何用规则、搜索和知识表示解决问题 & 智能是符号操作和问题求解过程 & 有助于分析推理、规划和可解释性 \\
联结主义传统 & 简单单元能否涌现复杂认知 & 智能来自分布式表示和连接权重更新 & 有助于理解深度学习和神经网络 \\
具身认知传统 & 身体和环境是否构成认知的一部分 & 智能嵌入感知、行动和情境互动 & 揭示 AI 缺乏身体经验的边界 \\
大模型研究传统 & 规模化数据和参数能否带来通用能力 & 智能表现可由大规模预训练和上下文学习增强 & 解释当前 AI 的能力跃迁与评测争议 \\
\bottomrule
\caption{人工智能与人类智能比较的研究脉络}
\label{tab:ai_human_literature}\\
\end{longtable}

表\ref{tab:ai_human_compare}概括了主要比较维度。两类智能系统的优势基础不同：AI 依赖可扩展计算、数据压缩和高速搜索；人类依赖具身经验、社会生活、主动目标和价值判断。

\begin{longtable}{p{2.6cm}p{4.0cm}p{4.0cm}p{3.2cm}}
\toprule
维度 & 人工智能的典型优势 & 人类智能的典型优势 & 关键边界 \\
\midrule
\endfirsthead
\toprule
维度 & 人工智能的典型优势 & 人类智能的典型优势 & 关键边界 \\
\midrule
\endhead
认知架构 & 可复制、可扩展、可在大规模硬件上并行部署 & 脑、身体、环境和社会关系共同构成认知系统 & AI 缺乏稳定具身经验 \\
信息处理 & 擅长高维统计模式识别和海量知识压缩 & 能把语言、行动、情感和目的结合成意义理解 & 统计相关不等于主体性理解 \\
学习方式 & 能从海量数据中学习复杂映射 & 低样本、主动探索、迁移和终身学习能力强 & AI 对数据分布和训练目标敏感 \\
推理能力 & 在形式化、可验证任务上进步显著 & 能结合常识、因果、反事实和责任判断 & benchmark 高分不等于开放世界推理稳定 \\
情感社会认知 & 能识别情绪线索并生成社交语言 & 有真实情绪、共情关系和伦理责任 & 模拟情绪不等于拥有情绪 \\
创造性 & 能快速生成大量候选方案并重组风格 & 能提出问题、选择价值和赋予作品社会意义 & 生成新组合不等于创造性主体 \\
效率与适应 & 数字任务速度快、可大规模服务用户 & 人脑低能耗，物理世界适应灵活 & 大模型训练成本高，机器人具身智能仍困难 \\
\bottomrule
\caption{人工智能与人类智能的多维度对比}
\label{tab:ai_human_compare}\\
\end{longtable}

\subsection{认知架构差异：符号、联结与具身系统}

人类智能是生物神经系统、身体和环境长期互动的结果。大脑与感官、运动系统、内分泌系统和社会环境共同构成认知系统。人类学习不仅发生在大脑内部，也发生在行动与反馈循环中。例如儿童学习“杯子”这一概念，除了看到大量图片，还会通过抓握、摔落、盛水、观察他人使用等多模态经验逐步形成概念。

人工智能的主流实现则以计算模型为核心。早期人工智能强调符号系统，用显式规则、逻辑推理和知识库表示智能。现代深度学习则主要基于联结主义思想，通过多层神经网络在大量数据中学习统计规律。大语言模型进一步将海量文本转化为高维参数空间中的语言知识，使模型可以生成连贯文本、回答问题和执行复杂指令。

二者的核心差异在于：人类认知是\textbf{具身的、情境化的、目标驱动的}；当前人工智能多是\textbf{数据驱动的、参数化的、任务约束的}。人工智能可以在没有身体经验的情况下通过文本学习“火是热的”，但它并没有被烫伤的感觉，也没有由身体经验形成的危险直觉。这种差异导致 AI 在语言层面可能表现得非常聪明，但在真实世界行动中仍容易出现常识缺失和情境误判。

\subsection{信息处理机制差异：统计相关与意义理解}

人工智能擅长从大规模数据中发现统计相关性。以深度神经网络为例，模型通过大量样本和梯度优化调整参数，使输入与输出之间形成复杂映射。大语言模型则通过预测下一个 token 学习语言分布，进而涌现出问答、摘要、翻译、代码生成和一定程度的推理能力。GPT-4 技术报告显示，GPT-4 在多项标准化考试中表现突出，例如 OpenAI 报告其在模拟律师考试中达到约第 90 百分位，在 MMLU 等知识测试中也显著优于早期模型\cite{openai2023gpt4}。

但是，统计相关还不能等同于真正理解。人类理解通常包含意图、因果、目的和语境。例如，人听到“他把钥匙放进冰箱里”会意识到这可能是反常行为，并联想到疲劳、粗心或特殊语境。AI 模型也可能给出类似解释，但这种解释主要来自训练数据中的模式匹配，缺少真实生活经验支撑。

这并不意味着人工智能完全不能形成抽象表示。深度模型中的中间层表征已经能够捕捉语义、句法、图像局部结构和跨模态关系。但这种表示通常缺乏人类式的主观体验和行动后果。人工智能可以形成\textbf{功能性意义}，即在任务中有效使用信息；至于它是否具有\textbf{体验性意义}和\textbf{主体性理解}，仍然是认知科学和人工智能哲学中的开放问题。

\subsection{学习与泛化：大数据学习与小样本学习}

人工智能的突出优势是大规模学习。现代模型可以在海量文本、图像、语音和代码上训练，从中学习人类个体无法直接接触的知识范围。例如 AlphaFold 系列模型通过学习蛋白质结构数据，在蛋白质结构预测领域取得突破；AlphaFold 3 进一步扩展到蛋白质与核酸、小分子、离子等多种生物分子的相互作用预测，被认为可能加速药物研发与基础生物学研究\cite{alphafold2024nature}。

然而，人类在小样本学习上更具优势。儿童看到几次新动物后就能在不同姿态、光照和背景中识别它；人类学习一个新词，也可以在少量语境中推断其意义并迁移使用。相比之下，人工智能通常需要大量训练样本，或者依赖预训练模型中已经存在的相似知识。虽然大语言模型具有上下文学习能力，可以通过少量示例完成新任务，但这种能力仍然依赖预训练中积累的大量统计经验。

泛化方面，AI 在\textbf{同分布泛化}上表现强，在\textbf{分布外泛化}上相对脆弱。一个图像模型在训练集相近的测试集上准确率很高，但面对风格变化、噪声、遮挡或对抗扰动时可能显著退化。人类也会犯错，但人类能够利用常识和环境反馈主动修正认知策略。当前人工智能通常缺乏这种稳定的自我校正机制。

\subsection{推理与因果理解：形式能力强，现实因果弱}

近年来，人工智能在数学推理、代码生成和逻辑问答方面进步显著。以 GPT-4 为例，技术报告显示其在多项标准化考试和复杂问答任务中明显优于早期模型；在医学考试等专业问答任务上，GPT-4 也被报告超过美国执业医师资格考试及格线二十余分\cite{openai2023gpt4,nori2023gpt4medical}。

但是，形式推理不等于完整的因果理解。人类因果推理通常包含反事实想象、干预判断和责任归因。例如，人类不仅能判断“雨后地面湿”，还知道“如果地面被洒水车洒过，即使没下雨也会湿”。当前大模型能够生成反事实推理文本，但其推理稳定性和可验证性仍有限。2025 年一篇关于机器类人智能基准的研究指出，许多 AI benchmark 缺乏人类验证标签，任务生态效度不足，不能简单用 benchmark 高分等同于人类式智能\cite{ying2025benchmarking}。

因此，可以说人工智能在\textbf{符号化、可文本化、可验证的推理任务}中进步很快，但在\textbf{真实世界因果推断、主动实验设计和复杂社会责任判断}方面仍不成熟。人类能将因果理解嵌入行动和价值判断中，而 AI 多数情况下仍停留在输入输出映射和语言解释层面。

\subsection{情感理解与社会智能：识别情绪不等于拥有情感}

人工智能可以识别人脸表情、语音情绪和文本情感倾向，也可以生成安慰、鼓励或共情风格的语言。这使 AI 在客服、心理健康辅助、教育陪伴和社交机器人中具有应用潜力。但是，AI 的情感理解本质上是对情感线索的模式识别和语言生成，并不等同于真正体验情绪。

人类情感具有身体基础和社会意义。愤怒、恐惧、悲伤、羞耻和爱不仅是认知判断，也伴随生理反应、记忆重构和行为倾向。人类共情还包含对他人处境的价值评价和责任感。AI 可以模拟共情语言，却不承担真实的情感风险和道德责任。因此，在高风险场景中，例如心理咨询、医疗沟通和法律建议，AI 更适合作为辅助工具，而不应完全替代人类专业判断。

社会智能方面，人类能够理解隐喻、礼貌、权力关系、群体规范和文化背景。AI 虽然能在大量文本中学习社会语言模式，但仍可能误解语境、刻板化群体或生成不合适建议。其根本原因在于 AI 缺乏长期社会生活经验，也缺乏真实参与社会关系的主体身份。

\subsection{能耗、效率与具身适应}

人脑约以几十瓦功率运行，却能完成感知、运动、语言、情绪、记忆和社会互动等复杂任务。大型 AI 模型的训练和推理则需要大量计算资源，尤其是大规模预训练阶段需要消耗巨量数据、算力和能源。AI 的优势是可以复制和部署到大量机器上，单个模型能力可以服务许多用户；但其训练成本、硬件依赖和环境影响也不可忽视。

具身适应方面，人类可以在复杂环境中连续感知、行动和学习。例如，一个人在陌生城市中可以通过问路、观察路标、避开障碍、理解他人表情来完成目标。当前 AI 在数字空间中表现很好，但进入物理世界后仍需要机器人、传感器、控制系统和安全机制配合。具身智能仍是未来人工智能的重要方向。

\subsection{优势与边界}

当前 AI 已经在一些明确任务中达到或超过普通人类水平：

（1）\textbf{大规模记忆与检索}：AI 能在海量文本中学习知识，快速生成相关内容。

（2）\textbf{高维模式识别}：在图像、语音、医学影像和结构预测等任务中，深度模型具有很强能力。

（3）\textbf{形式化任务求解}：数学、代码、标准化考试和逻辑问答中，强模型表现越来越接近专家水平。

（4）\textbf{重复性劳动自动化}：AI 能高效完成摘要、翻译、分类、格式转换和代码补全等任务。

（5）\textbf{科学辅助发现}：AlphaFold 等系统表明 AI 能帮助人类探索复杂科学空间，加速假设生成和实验设计。

（6）\textbf{生成与设计辅助}：AI 能快速生成文本、图像、代码和设计候选方案，适合头脑风暴和初稿生成。

这些优势来自可扩展计算、海量数据学习和高速搜索能力。不过，AI 仍难以匹敌人类的若干能力：

（1）\textbf{开放世界常识}：真实世界充满隐含规则、物理约束和社会语境，AI 容易在细节上犯低级错误。

（2）\textbf{低样本与终身学习}：人类可以在少量经验中学习并持续更新知识，而 AI 往往需要大量数据和重新训练。

（3）\textbf{因果与反事实推理}：AI 能生成推理文本，但不一定具有稳定的因果模型。

（4）\textbf{情绪和价值判断}：AI 能识别和模拟情绪，却没有真实情感体验和道德责任。

（5）\textbf{具身行动}：人类智能根植于身体与环境互动，AI 仍缺少稳定、安全、灵活的物理世界行动能力。

（6）\textbf{主体性和意义建构}：人类会追问“为什么做”，AI 多数情况下只优化“怎样做”。

在医学、教育和科学研究等场景中，AI 更适合承担检索、筛查、生成和辅助分析；最终判断、伦理沟通、责任承担和价值选择仍需要人类参与。

\subsection{批判性讨论：“超越人类”的语境与评测争议}

关于人工智能是否已经“超越人类”，最容易产生误解的地方在于评价标准不同。若评价标准是某个封闭任务的分数，例如分类准确率、棋类胜率、标准化考试成绩或代码题通过率，那么 AI 的确已经在许多任务中超过普通人，甚至接近或超过专家水平。但如果评价标准是开放世界中的综合适应能力、自主目标形成、伦理责任和长期社会生活能力，那么当前 AI 仍远未达到人类智能。

许多 benchmark 还存在生态效度问题。一个模型在选择题中表现良好，不一定能在真实实验室、医院、课堂或家庭环境中稳定行动。考试题通常有明确答案和固定上下文，现实任务则常常目标模糊、信息不完整、利益相关者复杂，并且错误会造成后果。AI 能力评估还会受到训练数据污染、提示工程、测试集分布和评价方式影响，因此应同时关注性能、鲁棒性、可解释性、安全性和责任边界。

\subsection{小结}

人工智能与人类智能不是同一把尺子上的简单高低关系。AI 在可数据化、可形式化、可评价的任务中进步很快，并在许多专业领域成为有效工具；人类则在具身经验、社会情境、价值判断和开放世界适应方面更有优势。因此，“超越人类”应限定在具体任务和评价指标中讨论。更现实的方向是人机协同：AI 承担大规模计算、检索、生成和优化任务，人类负责目标设定、价值判断、责任承担和复杂情境理解。

\newpage

\section{任务二：鸡尾酒会问题模拟：基于 TD-SpeakerBeam 的目标说话人提取系统}

\noindent\textbf{摘要：}
鸡尾酒会问题体现了人类听觉系统在多人声场中选择并持续跟踪目标说话人的能力。本文面向《认知与计算》课程任务，将该问题建模为目标说话人提取：系统输入双人混合语音和目标说话人的参考语音，输出目标说话人的估计语音。本文重点模拟选择性注意和目标声纹绑定，未显式建模空间声源定位；系统没有使用双耳时差/强度差、DOA、麦克风阵列或空间角度标签等空间线索。本文实现了基于 TD-SpeakerBeam 的目标说话人提取系统，并围绕数据构造、模型容量、教师--学生训练、频域前端和参考语音鲁棒性进行了系统评估。在 clean shared-80 数据设置上，可靠 SpeakerBeam 基线达到 10.32 dB test \sisdr{}，强教师模型达到 12.52 dB；中等容量学生模型达到 10.43 dB，fine-tune 蒸馏后达到 10.59 dB。进一步的 embedding-level 双参考语音聚合将学生模型提升到 10.80 dB，说明多条同说话人参考语音能够形成更稳定的目标声纹线索。enrollment sanity check 显示，正确参考语音下模型达到 10.59 dB，而 shuffled、interferer 和 zero enrollment 分别下降到 -7.32、-23.49 和 -6.38 dB，证明模型确实依据参考语音选择目标说话人。本文还记录了三个具有教学价值的诊断结果：早期数据构造中的 target/enrollment 泄漏会制造不可信高分；从零知识蒸馏在两种学生容量下均低于同架构监督训练；推理阶段直接加入未训练过的频谱加权前端会造成输入分布偏移。实验表明，可靠的目标说话人提取主干、严格的 enrollment 构造和针对参考语音的 sanity check 是构建鸡尾酒会语音系统的关键环节。

\noindent\textbf{关键词：}鸡尾酒会问题；目标说话人提取；TD-SpeakerBeam；听觉注意；语音分离；参考语音鲁棒性

\subsection{引言}

人在餐厅、教室和会议室中往往需要从多个同时说话的人声中跟踪某一个目标说话者。Cherry 的双耳语音识别实验将这一现象系统化为鸡尾酒会问题\cite{cherry1953}。从认知科学看，这一能力涉及选择性注意、听觉对象形成、特征绑定和资源调度等机制\cite{broadbent1958,bregman1990}。从计算角度看，它对应一个具有明确目标条件的语音分离问题：系统既要分离声音，也要判断“应该听谁”。

近年来，深度语音分离从 deep clustering\cite{hershey2016deep}、PIT\cite{yu2017pit}、Conv-TasNet\cite{luo2019convtasnet} 发展到 DPRNN\cite{luo2020dprnn}、SepFormer\cite{subakan2021sepformer} 和 TF-GridNet\cite{wang2023tfgridnet}，分离质量持续提升。然而，普通盲源分离通常假设输出源的排列可通过训练目标解决，并不直接模拟目标驱动的选择性听觉。目标说话人提取则更接近鸡尾酒会任务：参考语音提供目标身份线索，模型据此从混合语音中提取指定说话人。VoiceFilter、SpeakerBeam 和 SpEx+ 已经证明 speaker-conditioned extraction 是解决该任务的有效路径\cite{wang2019voicefilter,zmolikova2019speakerbeam,ge2020spexplus}。

本文围绕目标说话人提取系统的构建、评估和诊断展开，形成以下工作内容：
\begin{itemize}
    \item \textbf{目标说话人提取系统实现。}本文采用 TD-SpeakerBeam 构建 mixture + enrollment 到 target speech 的端到端系统，并在 80 位说话人的 clean shared 数据上完成训练与评估。
    \item \textbf{容量与教师--学生对照。}本文比较 strong teacher、small student 和 mid student，给出同架构监督训练与蒸馏训练的公平对照，分析容量、监督信号和教师输出之间的关系。
    \item \textbf{参考语音鲁棒性验证与多参考聚合。}本文评估 correct、shuffled、interferer、target、zero、short、noisy 等参考语音变体，验证模型是否真正执行目标说话人提取；进一步使用两条同说话人参考语音进行 embedding-level 聚合，将学生模型 \sisdr{} 提升到 10.80 dB。
    \item \textbf{实验诊断与失败案例。}本文记录数据构造失误、从零蒸馏和推理时频域前端三个诊断案例，并据此总结目标说话人提取实验中的数据审计、训练目标和训练/推理一致性要求。
\end{itemize}

\noindent\textbf{代码与脚本仓库：}
本项目的关键代码、训练与评估脚本、GUI 展示页面、真实 demo 音频和最终报告提交包已整理到 GitHub 仓库：
\url{https://github.com/CatNebulaaaa/CAFE-TSE}。其中，\texttt{src/cafe\_tse/} 保存目标说话人提取系统代码，\texttt{scripts/} 保存 TD-SpeakerBeam 训练、蒸馏、EGSP 探针、参考语音鲁棒性和 multi-enrollment 评估脚本，\texttt{gui/} 与 \texttt{demo\_audio/} 对应本文的交互式系统展示。

\subsection{认知背景与相关工作}

\subsubsection{听觉注意的计算映射}

人类鸡尾酒会能力通常包含三个层次。第一，选择性注意决定当前被优先处理的目标；第二，听觉场景分析把频率、时间、音色等线索组织成连续声源；第三，认知资源分配使系统在复杂场景中投入更多处理能力\cite{bregman1990}。在完整的人类听觉系统中，空间声源定位还会利用双耳时间差、双耳强度差和头相关传递函数等线索辅助注意转移；本文当前实现没有显式引入这些空间定位机制，而是将研究范围限定为单通道 mixture + enrollment 条件下的目标选择与声纹绑定。本文将这三点映射为计算系统：enrollment 提供目标线索，SpeakerBeam 的说话人适配层负责将目标身份注入分离网络，teacher/student 与容量对照实验用于分析计算资源和分离质量的关系。

\subsubsection{目标说话人提取}

目标说话人提取与盲源分离的区别在于，前者要求输出与指定参考语音一致的目标源。VoiceFilter 使用 speaker embedding 条件化语谱图掩码\cite{wang2019voicefilter}；SpeakerBeam 通过说话人感知神经网络提取目标源\cite{zmolikova2019speakerbeam}；SpEx+ 将参考语音条件引入时域提取网络\cite{ge2020spexplus}。这些工作说明，参考语音不仅是身份标签，也可以作为声学注意线索。本文采用 TD-SpeakerBeam 作为主要分离主干，并将错误参考语音、干扰说话人参考语音和零参考语音作为 sanity check。

\subsubsection{时频分离、条件调制与蒸馏}

时域分离网络避免显式相位估计，Conv-TasNet 等模型已经证明时域卷积结构能够有效处理语音分离\cite{luo2019convtasnet}。TF-GridNet 等时频模型则强调 full-band 与 sub-band 结构建模\cite{wang2023tfgridnet}。本文主要采用 TD-SpeakerBeam，是因为其目标说话人适配机制与 enrollment-based selective attention 高度一致。训练方面，知识蒸馏通常希望用教师模型的输出作为软目标，使小模型继承大模型的行为\cite{hinton2015distill}；本文在 clean target 监督充分的设置下重新检验了这一策略，并将其与同架构监督学生模型公平比较。

\subsection{方法}

\subsubsection{任务定义}

给定混合语音 $x\in\mathbb{R}^{T}$、目标说话人参考语音 $e\in\mathbb{R}^{T_e}$ 和目标干净语音 $s\in\mathbb{R}^{T}$，目标说话人提取学习函数
\begin{equation}
    \hat{s}=f_{\theta}(x,e),
\end{equation}
使估计语音 $\hat{s}$ 与 $s$ 尽可能一致，同时抑制非目标说话人。本文采用 \sisdr{} loss 训练监督模型，并使用 BSS Eval 风格指标评价分离质量\cite{vincent2006bss,leroux2019sdr}。在蒸馏实验中，学生模型同时接收真实目标语音和教师模型输出作为训练信号。

\subsubsection{系统结构}

图 \ref{fig:arch} 展示了系统结构。系统由五个环节共同构成：第一，数据构造环节从干净单说话人语音中合成双人 mixture，并为目标说话人选择同说话人的另一句 enrollment；第二，参考语音分支把 enrollment 编码为目标说话人 embedding；第三，TD-SpeakerBeam 主干用该 embedding 调制时域分离网络，输出目标波形；第四，训练环节包含监督学习、教师--学生对照和 fine-tune 蒸馏；第五，评估环节包含标准客观指标、错误参考语音 sanity check、短/噪参考鲁棒性、频域前端探针和多参考语音聚合。分离主干解决“如何分离”，参考、数据和评估环节负责明确“听谁、数据是否可信、结果如何验证”。

\begin{figure}[H]
\centering
\begin{tikzpicture}[
    scale=1.10, transform shape,
    font=\small,
    core/.style={rectangle, rounded corners=4pt, draw=tseBlueLine, fill=tseBlue, minimum width=3.0cm, minimum height=0.88cm, align=center},
    input/.style={rectangle, rounded corners=4pt, draw=tseAmberLine, fill=tseAmber, minimum width=2.8cm, minimum height=0.82cm, align=center},
    verify/.style={rectangle, rounded corners=4pt, draw=tseGreenLine, fill=tseGreen, minimum width=2.95cm, minimum height=0.78cm, align=center},
    note/.style={rectangle, rounded corners=4pt, draw=tseGrayLine, fill=tseGray, minimum width=2.95cm, minimum height=0.72cm, align=center},
    arr/.style={-{Latex[length=2.3mm]}, thick, draw=black!72},
    soft/.style={-{Latex[length=2.0mm]}, thick, dashed, draw=black!55}
]
\node[input] (mix) at (-5.4,0.7) {混合语音 $x$\\目标 + 干扰};
\node[core] (enc) at (-2.05,0.7) {时域编码器\\learned filterbank};
\node[core] (mask) at (1.55,0.7) {TD-SpeakerBeam\\说话人条件掩码};
\node[input] (out) at (5.15,0.7) {目标输出\\$\hat{s}$};
\node[input] (enroll) at (-5.4,-1.15) {参考语音 $e$\\同人不同句};
\node[core] (embed) at (-2.05,-1.15) {Auxiliary 分支\\speaker embedding};
\node[verify] (pool) at (1.55,-1.15) {多参考聚合\\embedding pooling};
\node[verify] (probe) at (5.15,-1.15) {参考语音探针\\wrong/short/noisy};
\node[note] (train) at (-3.7,-2.75) {训练对照\\teacher/student/distill};
\node[note] (eval) at (3.45,-2.75) {指标评估\\SI-SDR/SDR/SIR/SAR/STOI/PESQ};

\draw[arr] (mix) -- (enc);
\draw[arr] (enc) -- (mask);
\draw[arr] (mask) -- (out);
\draw[arr] (enroll) -- (embed);
\draw[arr] (embed) -- (mask.south west);
\draw[arr] (embed.east) -- (pool.west);
\draw[arr] (pool.north) -- (mask.south);
\draw[soft] (probe.north west) -- (mask.south east);
\draw[soft] (train) -| (mask.south);
\draw[soft] (out.south) |- (eval.east);
\draw[soft] (probe.south) |- (eval.east);

\begin{scope}[on background layer]
\node[draw=tseBlueLine!45, fill=tseBlueLine!4, rounded corners=6pt, fit=(mix)(enc)(mask)(out), inner xsep=8pt, inner ysep=7pt] {};
\node[draw=tseGreenLine!45, fill=tseGreenLine!4, rounded corners=6pt, fit=(enroll)(embed)(pool)(probe), inner xsep=8pt, inner ysep=7pt] {};
\end{scope}
\end{tikzpicture}
\caption{目标说话人提取系统结构。实线为分离主路径，虚线为训练、诊断和多参考聚合环节。}
\label{fig:arch}
\end{figure}

\subsubsection{数据构造、训练与评估流程}

图 \ref{fig:data_train_flow} 给出本文从数据构造到模型评估的完整流程。该流程把数据审计单独列出，是因为目标说话人提取对 enrollment 构造非常敏感：参考语音必须来自同一说话人的另一句语音，不能与目标语音泄漏为同一路径。训练完成后，本文同时进行标准 test 评估和 enrollment sanity check，以区分真正的目标说话人提取与普通语音增强。

\begin{figure}[H]
\centering
\begin{tikzpicture}[
    scale=1.05, transform shape,
    font=\small,
    flow/.style={rectangle, rounded corners=4pt, draw=tseBlueLine, fill=tseBlue, minimum width=2.65cm, minimum height=0.78cm, align=center},
    audit/.style={rectangle, rounded corners=4pt, draw=tseGreenLine, fill=tseGreen, minimum width=2.85cm, minimum height=0.78cm, align=center},
    result/.style={rectangle, rounded corners=4pt, draw=tseAmberLine, fill=tseAmber, minimum width=3.05cm, minimum height=0.78cm, align=center},
    arr/.style={-{Latex[length=2.2mm]}, thick, draw=black!72}
]
\node[flow] (speech) at (-4.8,1.65) {LibriSpeech\\clean speech};
\node[flow] (sample) at (0,1.65) {shared-80 split\\train/valid/test};
\node[flow] (mix) at (4.8,1.65) {mixture 构造\\target + interferer};
\node[audit] (audit) at (-2.4,0.15) {Manifest 审计\\speaker/path/disjoint};
\node[flow] (train) at (2.4,0.15) {模型训练\\supervised/distill};

\node[result] (metric) at (-4.8,-1.45) {标准测试\\SI-SDR/SDR/SIR/SAR};
\node[result] (sanity) at (0,-1.45) {参考条件探针\\wrong/short/noisy};
\node[result] (egsp) at (4.8,-1.45) {增强模块探针\\EGSP/pooling};
\node[audit] (conclusion) at (0,-3.05) {统一结论\\质量、鲁棒性与失败诊断};

\draw[arr] (speech) -- (sample);
\draw[arr] (sample) -- (mix);
\draw[arr] (mix.south) |- (audit.east);
\draw[arr] (audit) -- (train);
\draw[arr] (train.south) |- (metric.east);
\draw[arr] (train.south) -- (sanity.north);
\draw[arr] (train.south) |- (egsp.west);
\draw[arr] (metric.south) |- (conclusion.west);
\draw[arr] (sanity) -- (conclusion);
\draw[arr] (egsp.south) |- (conclusion.east);
\end{tikzpicture}
\caption{数据构造、训练与评估流程。数据审计和参考语音探针用于保证模型确实学习目标说话人条件。}
\label{fig:data_train_flow}
\end{figure}

\subsubsection{监督训练与教师--学生训练}

监督学生模型直接优化估计语音与干净目标语音之间的 \sisdr{} loss：
\begin{equation}
    \mathcal{L}_{target}=\mathcal{L}_{\sisdr{}}(\hat{s},s).
\end{equation}
蒸馏模型在此基础上加入教师输出 $\hat{s}_{T}$：
\begin{equation}
    \mathcal{L}=\mathcal{L}_{target}+\lambda\mathcal{L}_{\sisdr{}}(\hat{s},\hat{s}_{T}).
\end{equation}
本文比较两种蒸馏方式：一种从随机初始化开始训练，另一种从同架构监督学生 checkpoint 出发，用较小教师权重和较低学习率 fine-tune。这样可以区分“教师信号本身的帮助”和“已有监督模型上的弱正则化”。

\begin{figure}[H]
\centering
\begin{tikzpicture}[
    scale=1.02, transform shape,
    font=\small,
    plain/.style={rectangle, rounded corners=4pt, draw=tseGrayLine, fill=tseGray, minimum width=2.8cm, minimum height=0.8cm, align=center},
    stu/.style={rectangle, rounded corners=4pt, draw=tseBlueLine, fill=tseBlue, minimum width=3.05cm, minimum height=0.8cm, align=center},
    tea/.style={rectangle, rounded corners=4pt, draw=tseAmberLine, fill=tseAmber, minimum width=3.05cm, minimum height=0.8cm, align=center},
    bad/.style={rectangle, rounded corners=4pt, draw=tseRedLine, fill=tseRed, minimum width=3.05cm, minimum height=0.8cm, align=center},
    resultbox/.style={rectangle, rounded corners=4pt, draw=tseGreenLine, fill=tseGreen, minimum width=3.2cm, minimum height=0.8cm, align=center},
    arr/.style={-{Latex[length=2.2mm]}, thick, draw=black!72},
    faint/.style={-{Latex[length=2.0mm]}, thick, dashed, draw=black!55}
]
\node[plain] (input) at (-5.4,0) {mixture + enrollment};
\node[stu] (supervised) at (-2.0,1.05) {监督学生\\clean target};
\node[tea] (teacher) at (-2.0,-1.05) {Strong teacher\\teacher waveform};
\node[bad] (scratch) at (1.75,1.05) {从零蒸馏\\低于监督基线};
\node[stu] (ft) at (1.75,-1.05) {Fine-tune 蒸馏\\小幅正则收益};
\node[resultbox] (compare) at (5.25,0) {同容量测试\\质量与干扰抑制};
\node[plain] (takeaway) at (1.75,-2.65) {结论\\蒸馏不能替代可靠监督训练};

\draw[arr] (input) -- (supervised);
\draw[arr] (input) -- (teacher);
\draw[arr] (supervised) -- (scratch);
\draw[faint] (teacher) -- (scratch);
\draw[arr] (supervised) -- (ft);
\draw[faint] (teacher) -- (ft);
\draw[arr] (scratch) -- (compare);
\draw[arr] (ft) -- (compare);
\draw[arr] (compare.south) |- (takeaway.east);
\end{tikzpicture}
\caption{知识蒸馏实验设计。本文将同架构监督学生、从零蒸馏和 fine-tune 蒸馏放在同一测试集上比较。}
\label{fig:distill_flow}
\end{figure}

\subsubsection{频域前端与参考语音探针}

为分析参考语音频谱画像能否作为显式前端，本文实现了推理时 EGSP 探针。给定 enrollment 幅度谱 $E$，先计算长期频谱画像
\begin{equation}
    p_f=\frac{1}{N_e}\sum_{t=1}^{N_e}E_{f,t},
\end{equation}
再对 mixture 的 STFT 幅度施加频带权重：
\begin{equation}
    w_f=\mathrm{clip}(1+\alpha(p_f/\mu(p)-1),w_{\min},w_{\max}).
\end{equation}
该探针不更新模型参数，用于观察未训练频域前端是否能直接提升已训练 SpeakerBeam。参考语音探针则替换 enrollment 输入，包括 correct、shuffled、interferer、target-as-enrollment、zero、short 和 noisy variants。

\begin{figure}[H]
\centering
\begin{tikzpicture}[
    scale=1.12, transform shape,
    font=\small,
    good/.style={rectangle, rounded corners=4pt, draw=tseGreenLine, fill=tseGreen, minimum width=3.05cm, minimum height=0.78cm, align=center},
    bad/.style={rectangle, rounded corners=4pt, draw=tseRedLine, fill=tseRed, minimum width=3.05cm, minimum height=0.78cm, align=center},
    mid/.style={rectangle, rounded corners=4pt, draw=tseAmberLine, fill=tseAmber, minimum width=3.05cm, minimum height=0.78cm, align=center},
    core/.style={rectangle, rounded corners=4pt, draw=tseBlueLine, fill=tseBlue, minimum width=3.35cm, minimum height=0.86cm, align=center},
    resultbox/.style={rectangle, rounded corners=4pt, draw=tseGrayLine, fill=tseGray, minimum width=3.25cm, minimum height=0.8cm, align=center},
    arr/.style={-{Latex[length=2.2mm]}, thick, draw=black!72}
]
\node[resultbox] (mix) at (0,1.85) {固定 mixture\\同一段混合语音};
\node[good] (correct) at (-4.7,1.35) {正确身份\\correct / target};
\node[mid] (degrade) at (-4.6,0) {退化参考\\short / noisy / zero};
\node[bad] (wrong) at (-4.7,-1.35) {错误身份\\shuffled / interferer};
\node[core] (model) at (0,0) {固定 TD-SpeakerBeam\\只替换 enrollment};
\node[resultbox] (metric) at (4.6,0) {输出比较\\SI-SDR、SIR、STOI、PESQ};
\node[good] (pass) at (4.6,1.55) {身份验证通过\\正确参考高分};
\node[bad] (fail) at (4.6,-1.55) {身份验证失败\\错误参考崩溃};

\draw[arr] (mix) -- (model);
\draw[arr] (correct) -- (model);
\draw[arr] (degrade) -- (model);
\draw[arr] (wrong) -- (model);
\draw[arr] (model) -- (metric);
\draw[arr] (metric) -- (pass);
\draw[arr] (metric) -- (fail);
\end{tikzpicture}
\caption{Enrollment sanity check 设计。固定 mixture 和模型，仅替换参考语音以验证目标身份条件是否生效。}
\label{fig:enroll_probe}
\end{figure}

\subsection{实验设置}

本文使用基于 LibriSpeech clean 语音构造的 shared-clean80 目标说话人提取数据。LibriSpeech 是公开英文有声书语音数据集，原始语音为单说话人干净朗读；本文不直接使用其原始识别标签，而是从中抽取 80 位说话人的干净语音片段，按目标说话人提取任务重新生成 manifest。每条样本包含四个音频路径：目标说话人语音 target、另一位说话人的 interferer、二者相加得到的 mixture，以及来自目标说话人但与 target 不同句的 enrollment。这样构造后，任务输入为 mixture 和 enrollment，输出为 target，符合鸡尾酒会场景中“多人同时说话，但只听某个目标说话人”的设定。

数据规模为 8000 条训练样本、800 条验证样本和 800 条测试样本。训练与评估均采用 8 kHz 采样率和 4 s 片段；训练、验证和测试 manifest 分开生成，最终结果均要求 target 与 enrollment 为同一说话人的不同语音片段，避免把目标语音本身泄漏为参考语音。早期实验曾出现 target/enrollment 路径重合和 speaker id 解析错误，因此最终流程中单独加入 manifest 审计，检查 speaker id、路径是否存在、target/enrollment 是否 disjoint，以及 interferer 是否来自非目标说话人。

评价指标包括 \sisdr{}、\sisdr{}i、SDR、SIR、SAR、STOI 和 PESQ。\sisdr{} 是尺度不变信号失真比，会在计算前允许估计信号按最优比例缩放，因此更关注波形形状是否接近目标；传统 SDR 则在给定参考源分解的条件下衡量总体失真。本文的 mixture、target 和估计语音在评估前经过统一长度和归一化处理，且是单目标提取任务，因此 \sisdr{} 与 SDR 经常非常接近，但二者定义并不相同。SIR 衡量干扰说话人抑制程度，SAR 衡量算法伪影和失真，STOI 近似反映可懂度，PESQ 近似反映感知语音质量。实验分为七组：主干模型对比、学生容量对照、蒸馏消融、不同混合 SNR 与说话者数量压力测试、EGSP 推理探针、enrollment sanity check 和 multi-enrollment 聚合。其中，混合 SNR 实验在原 clean two-speaker test set 上额外加入 babble-like 噪声，设置 20、10、5、0 dB 四个 SNR；说话者数量实验保持目标说话人与 enrollment 不变，将第三位非目标说话人加入 mixture，形成 3-speaker clean 压力测试。

\subsection{实验结果}

\subsubsection{主结果}

表 \ref{tab:main} 给出主结果。早期自写 CAFE-TSE 原型在严格目标说话人提取设置下只有约 1 dB \sisdr{}，说明其条件建模和数据构造尚不足以支撑最终系统。切换到 TD-SpeakerBeam 后，系统在 shared-clean80 数据上稳定达到 10 dB 以上；更大容量 strong teacher 进一步达到 12.52 dB。中等容量学生模型在保持较低容量的同时达到 10.43 dB，成为后续学生实验的主要基线。表 \ref{tab:final-full-metrics} 进一步给出最终学生相关变体的完整多指标结果。

\begin{table}[H]
\centering
\small
\begin{tabular}{@{}p{4.7cm}rrrp{2.4cm}@{}}
\toprule
系统 & \sisdr{}$\uparrow$ & SDR$\uparrow$ & SIR$\uparrow$ & 角色 \\
\midrule
Self-written CAFE-TSE & 1.06 & 1.06 & 4.56 & 失败原型 \\
Open SpeakerBeam shared-clean80 mid & 10.32 & 10.34 & 22.23 & 可靠基线 \\
Strong teacher & \textbf{12.52} & \textbf{12.54} & \textbf{27.05} & 容量上界 \\
Small supervised student & 8.16 & 8.16 & 17.30 & 小容量学生 \\
Mid supervised student & 10.43 & 10.43 & 21.82 & 主要学生基线 \\
Mid fine-tune distill & 10.59 & 10.59 & 22.90 & 单参考学生 \\
Mid distill + 2-enrollment embedding pooling & 10.80 & 10.80 & 23.22 & 最终学生变体 \\
\bottomrule
\end{tabular}
\normalsize
\caption{主要系统结果。SpeakerBeam 主干显著高于早期自写原型，strong teacher 给出当前质量上界，mid student 是本文主要学生基线。}
\label{tab:main}
\end{table}

\begin{table}[H]
\centering
\small
\begin{tabular}{lrrrrrrr}
\toprule
系统 & \sisdr{} & \sisdr{}i & SDR & SIR & SAR & STOI & PESQ \\
\midrule
单参考 mid fine-tune & 10.586 & 10.581 & 10.591 & 22.899 & 11.398 & 0.544 & 1.308 \\
2-enroll embedding pooling & \textbf{10.804} & \textbf{10.800} & \textbf{10.799} & \textbf{23.224} & \textbf{11.517} & \textbf{0.546} & \textbf{1.313} \\
4-enroll embedding pooling & 10.684 & 10.679 & 10.681 & 23.068 & 11.488 & 0.545 & 1.311 \\
8-enroll embedding pooling & 10.662 & 10.657 & 10.660 & 23.083 & 11.472 & 0.545 & 1.310 \\
\bottomrule
\end{tabular}
\normalsize
\caption{最终学生变体的完整多指标结果。多参考语音在 embedding 层聚合后，$k=2$ 同时改善失真、干扰抑制、伪影和感知指标。}
\label{tab:final-full-metrics}
\end{table}

\subsubsection{学生容量分析}

表 \ref{tab:capacity} 比较 small 和 mid 两种学生容量。small student 使用 128 通道、4 blocks、2 repeats；mid student 使用 256 通道、6 blocks、2 repeats。容量提升将 \sisdr{} 从 8.16 dB 提高到 10.43 dB，并显著提高 SIR。这说明在该 clean TSE 设置中，学生容量对分离质量具有直接影响。

\begin{table}[H]
\centering
\begin{tabular}{lrrrr}
\toprule
学生模型 & n\_filters/hid & blocks & Test \sisdr{}$\uparrow$ & Test SIR$\uparrow$ \\
\midrule
Small supervised & 128 & 4 & 8.16 & 17.30 \\
Mid supervised & 256 & 6 & \textbf{10.43} & \textbf{21.82} \\
\bottomrule
\end{tabular}
\caption{学生容量对照。增大学生容量带来 2.27 dB \sisdr{} 提升。}
\label{tab:capacity}
\end{table}

\subsubsection{知识蒸馏消融}

表 \ref{tab:distill-failure} 给出两种容量下的蒸馏实验。from-scratch distillation 在 small 和 mid 学生上均低于同架构监督训练；从监督 checkpoint 初始化并以小权重使用教师输出进行 fine-tune 时，small 和 mid 学生分别获得 0.17 dB 和 0.16 dB 的小幅提升。

\begin{table}[H]
\centering
\begin{tabular}{lrrrl}
\toprule
模型 & \sisdr{}$\uparrow$ & SDR$\uparrow$ & SIR$\uparrow$ & 结论 \\
\midrule
Strong teacher & 12.52 & 12.54 & 27.05 & 教师上界 \\
Small supervised & 8.16 & 8.16 & 17.30 & 小容量监督基线 \\
Small from-scratch distill & 8.02 & 8.02 & 16.76 & 低于监督训练 \\
Small fine-tune distill & 8.33 & 8.34 & 17.38 & 小幅提升 \\
Mid supervised & 10.43 & 10.43 & 21.82 & 中等容量监督基线 \\
Mid from-scratch distill & 10.26 & 10.27 & 21.80 & 低于监督训练 \\
Mid fine-tune distill & 10.59 & 10.59 & 22.90 & 小幅提升 \\
\bottomrule
\end{tabular}
\caption{强教师蒸馏消融。两种学生容量下，从零蒸馏均未超过同架构监督训练；fine-tune 蒸馏只有约 0.16--0.17 dB 的边际收益。}
\label{tab:distill-failure}
\end{table}

该结果与语音波形回归任务的监督结构一致：训练集中已经提供干净目标语音 $s$，教师输出 $\hat{s}_{T}$ 仍包含残余失真。较大的 teacher loss 会让学生同时靠近 ground truth 和 teacher residual。fine-tune 蒸馏可作为弱正则使用，但容量提升和可靠主干选择对结果影响更大。

\subsubsection{频域 EGSP 探针}

表 \ref{tab:egsp-probe} 给出推理时 EGSP strength sweep。该实验只在推理阶段对 mixture 做 STFT 频带加权，不重新训练模型。结果显示，所有非零 strength 均低于 0.00 基线，说明未训练的频域前端会改变 TD-SpeakerBeam 的输入分布。随后本文又尝试在训练阶段加入 EGSP fine-tune：strength=0.02 得到 10.557 dB \sisdr{}，strength=0.05 得到 10.577 dB \sisdr{}，仍未超过单参考 mid fine-tune 的 10.586 dB。该结果说明频域先验并非不可用，但当前简单乘性加权不足以稳定提升时域 SpeakerBeam。

\begin{table}[H]
\centering
\begin{tabular}{@{}S[table-format=1.2]S[table-format=2.3]S[table-format=2.3]S[table-format=2.3]S[table-format=2.3]S[table-format=1.3]@{}}
\toprule
{\makebox[2.7cm][c]{EGSP strength}} & {\sisdr{}$\uparrow$} & {SDR$\uparrow$} & {SIR$\uparrow$} & {SAR$\uparrow$} & {STOI$\uparrow$} \\
\midrule
0.00 & {\bfseries 10.586} & {\bfseries 10.591} & {\bfseries 22.899} & {\bfseries 11.398} & {\bfseries 0.544} \\
0.02 & 10.464 & 10.469 & 22.881 & 11.273 & 0.542 \\
0.05 & 10.348 & 10.353 & 22.820 & 11.156 & 0.540 \\
0.10 & 10.288 & 10.293 & 22.785 & 11.092 & 0.538 \\
0.20 & 10.186 & 10.190 & 22.749 & 10.986 & 0.537 \\
0.30 & 10.079 & 10.082 & 22.766 & 10.875 & 0.536 \\
\bottomrule
\end{tabular}
\caption{EGSP 推理探针。未经过训练的频域加权在所有 strength 下均降低分离指标。}
\label{tab:egsp-probe}
\end{table}

该探针给出一个明确的实现边界：若要继续使用 enrollment-guided frequency weighting，应在训练阶段端到端加入，使模型适应加权后的波形分布；仅把它作为推理后处理会损害已经训练好的时域分离器。

\subsubsection{混合 SNR 与说话者数量变化}

表 \ref{tab:stress} 给出不同混合 SNR 和说话者数量下的补充实验。该实验固定 mid fine-tune distill 模型，不重新训练；因此它评估的是最终系统面对更复杂鸡尾酒会条件时的泛化能力。clean 2-speaker 条件下模型达到 10.586 dB \sisdr{}。当 mixture 额外加入 babble-like 噪声时，20 dB 条件仍保持 9.770 dB；10 dB、5 dB 和 0 dB 条件分别下降到 5.941、3.358 和 0.906 dB，说明未进行噪声增强训练的模型对强背景噪声较敏感。3-speaker clean 条件下 \sisdr{} 为 2.674 dB，但 \sisdr{}i 仍有 7.182 dB，说明模型仍能利用 enrollment 抽取目标说话人，只是额外干扰源显著增加了抑制难度。

\begin{table}[H]
\centering
\small
\begin{tabular}{lrrrrrr}
\toprule
测试条件 & \sisdr{}$\uparrow$ & \sisdr{}i$\uparrow$ & SDR$\uparrow$ & SIR$\uparrow$ & STOI$\uparrow$ & PESQ$\uparrow$ \\
\midrule
2-speaker clean & \textbf{10.586} & \textbf{10.581} & \textbf{10.591} & \textbf{22.899} & \textbf{0.544} & \textbf{1.309} \\
2-speaker + babble 20 dB & 9.770 & 9.873 & 9.775 & 21.987 & 0.535 & 1.306 \\
2-speaker + babble 10 dB & 5.941 & 6.884 & 5.944 & 16.578 & 0.520 & 1.305 \\
2-speaker + babble 5 dB & 3.358 & 5.782 & 3.359 & 13.362 & 0.509 & 1.278 \\
2-speaker + babble 0 dB & 0.906 & 6.075 & 0.906 & 10.187 & 0.495 & 1.274 \\
3-speaker clean & 2.674 & 7.182 & 2.676 & 9.638 & 0.504 & 1.279 \\
\bottomrule
\end{tabular}
\normalsize
\caption{不同混合 SNR 与说话者数量下的压力测试。强噪声和第三说话者会显著降低绝对分离质量，但 \sisdr{}i 显示系统仍能提供目标说话人相关增益。}
\label{tab:stress}
\end{table}

\subsubsection{Enrollment sanity check}

表 \ref{tab:enroll} 给出参考语音变体实验。correct enrollment 下模型达到 10.586 dB；shuffled、interferer 和 zero enrollment 显著下降，interferer enrollment 甚至降至 -23.493 dB。这说明模型的输出由参考语音身份强烈控制，具备目标说话人选择能力。短参考语音和带噪参考语音给出渐进退化：1s enrollment 仍有 9.191 dB，2s enrollment 接近完整参考语音；10 dB noisy enrollment 基本可用，5 dB noisy enrollment 出现明显下降。

\begin{table}[H]
\centering
\begin{tabular}{lrrrrrr}
\toprule
Enrollment & \sisdr{}$\uparrow$ & SDR$\uparrow$ & SIR$\uparrow$ & SAR$\uparrow$ & STOI$\uparrow$ & PESQ$\uparrow$ \\
\midrule
correct & 10.586 & 10.591 & 22.899 & 11.398 & 0.544 & 1.308 \\
shuffled & -7.318 & -7.303 & -0.261 & 6.255 & 0.411 & 1.274 \\
interferer & -23.493 & -23.462 & -22.975 & 11.487 & 0.311 & 1.253 \\
target & \textbf{10.977} & \textbf{10.972} & \textbf{23.349} & \textbf{11.636} & \textbf{0.547} & 1.312 \\
zero & -6.381 & -6.318 & -0.414 & 7.292 & 0.429 & 1.276 \\
short 1s & 9.191 & 9.203 & 20.931 & 10.597 & 0.535 & 1.308 \\
short 2s & 10.366 & 10.371 & 22.664 & 11.241 & 0.543 & \textbf{1.311} \\
noisy 10 dB & 10.085 & 10.097 & 21.857 & 11.127 & 0.542 & 1.308 \\
noisy 5 dB & 7.318 & 7.336 & 17.782 & 9.719 & 0.525 & 1.301 \\
\bottomrule
\end{tabular}
\caption{参考语音 sanity check 与鲁棒性分析。错误参考语音会使目标身份选择崩溃，短参考和轻噪参考只造成渐进退化。}
\label{tab:enroll}
\end{table}

\subsubsection{Multi-enrollment 聚合}

表 \ref{tab:multi-enroll} 给出多参考语音聚合实验。该实验固定 mid fine-tune distill 模型，不改变训练参数；对每条测试样本，除原 enrollment 外，从训练集同一说话人中选取额外参考语音，先分别经过 SpeakerBeam auxiliary 分支得到 enrollment embedding，再在 embedding 层求平均并送入 masker。$k=2$ 时 \sisdr{} 从 10.586 dB 提升到 10.804 dB，同时 SDR、SIR、STOI 和 PESQ 也略有提高。$k=4$ 和 $k=8$ 没有继续提升，说明参考语音数量并非越多越好：少量额外参考可以稳定目标声纹，过多参考则可能引入录音内容、音量和发音状态差异。

\begin{table}[H]
\centering
\begin{tabular}{rrrrrr}
\toprule
参考条数 $k$ & \sisdr{}$\uparrow$ & SDR$\uparrow$ & SIR$\uparrow$ & STOI$\uparrow$ & PESQ$\uparrow$ \\
\midrule
1 & 10.586 & 10.591 & \textbf{22.899} & 0.544 & 1.309 \\
2 & \textbf{10.804} & \textbf{10.799} & \textbf{23.224} & \textbf{0.546} & \textbf{1.313} \\
4 & 10.684 & 10.681 & 23.068 & 0.545 & 1.311 \\
8 & 10.662 & 10.660 & 23.083 & 0.545 & 1.310 \\
\bottomrule
\end{tabular}
\caption{Multi-enrollment embedding-level 聚合。两条同说话人参考语音带来最稳定的目标声纹线索。}
\label{tab:multi-enroll}
\end{table}

\subsection{失败案例与诊断反思}

失败案例是本文实验反思中最关键的部分，因为它说明语音分离系统不能只依赖单个高分表格。目标说话人提取同时依赖数据构造、参考语音条件、模型主干和评价指标；任一环节出错，都可能得到看似合理但实际不可解释的结果。表 \ref{tab:failure} 总结了三个主要诊断案例，随后分别讨论其原因与修正经验。

\begin{table}[H]
\centering
\begin{tabular}{p{3.2cm}p{5.0cm}p{5.2cm}}
\toprule
案例 & 现象 & 处理与经验 \\
\midrule
自写原型与数据构造失误 & 早期 CAFE-TSE 约 1 dB；参考锚点在泄漏设置下产生虚高；\texttt{s2} speaker 解析曾导致目标说话人错标 & 将自写原型放入失败案例，修复 target/enrollment disjoint 和 speaker id 解析，主结果采用 TD-SpeakerBeam \\
从零知识蒸馏 & small 与 mid 学生从零蒸馏均低于监督训练 & 在 clean target 监督充分时，教师波形只适合作为低权重 fine-tune 正则 \\
未训练 EGSP 前端 & 推理时加入 EGSP 后，\sisdr{} 从 10.586 下降到 10.079--10.464 & 频域前端若用于性能提升，应训练时端到端集成，不能作为后处理硬插入 \\
\bottomrule
\end{tabular}
\caption{失败案例与诊断结论。}
\label{tab:failure}
\end{table}

\textbf{第一，数据泄漏会制造虚假的“创新”。}早期实验曾尝试把参考语音频谱或锚点直接用于输出约束，在部分设置下得到过较高分数。但进一步检查发现，若 enrollment 与 target 路径重合，或者 speaker id 解析把干扰说话人误标为目标说话人，模型可能利用近似目标本身的信息，而不是学习真正的目标说话人提取。修复方式是强制 target/enrollment disjoint，并在 manifest 审计中检查 speaker id、路径、源说话人和 interferer 关系。这个案例说明，目标说话人提取比普通去噪更容易受到参考语音泄漏影响，任何声称“参考语音增强有效”的模块都必须先通过 wrong enrollment 和 interferer enrollment 检查。

\textbf{第二，自写主干失败暴露了条件建模难度。}自写 CAFE-TSE 原型包含 EGSP、条件融合、轻量 TF-GridNet 和动态推理等多个模块，但在严格 disjoint enrollment 下只有约 1 dB \sisdr{}。这说明目标说话人提取的核心难点在于稳定地把“目标是谁”注入分离网络。若主干本身不能可靠建模语音时序和说话人条件，再多外围模块也难以补救。因此最终报告采用 TD-SpeakerBeam 作为可靠主干，把自写原型保留为失败案例和系统设计探索。

\textbf{第三，蒸馏不是天然有效。}在图像分类中，教师模型的软标签常常能提供类别间关系；但本文的波形回归任务已经有干净 target 监督，教师输出仍可能包含残余干扰和伪影。从零蒸馏同时拟合 clean target 与 teacher waveform，反而低于同架构监督训练。只有从监督学生 checkpoint 出发、用较小教师权重 fine-tune 时，才获得约 0.16--0.17 dB 的边际收益。因此本文不把蒸馏包装成主要创新，而把它作为一个有条件的小幅正则化手段。

\textbf{第四，频域先验必须与训练分布一致。}EGSP 的动机来自人类听觉注意：参考语音可以提示目标说话人的长期频谱特征。但直接在推理阶段给 mixture 做频带加权，会让已训练的 TD-SpeakerBeam 接收到分布外输入，导致 \sisdr{}、SAR 和 STOI 同时下降。训练阶段 EGSP fine-tune 也未超过单参考基线，说明简单乘性频带权重仍不够。更合理的后续方向是把频域画像作为可学习条件或辅助损失，而不是作为固定前端硬插入。

这些失败案例共同指向一个方法论结论：鸡尾酒会系统的可靠性来自“主干有效、数据无泄漏、参考语音可验证、指标可解释”四个条件。相比单纯堆叠模块，严格的对照实验和负结果解释更能体现系统是否真正完成了目标说话人提取。

\subsection{认知机制讨论}

\subsubsection{从选择性注意到目标说话人条件}

人类听觉注意会把目标相关线索置于优先地位。本文中的 enrollment 就是计算系统中的注意提示。表 \ref{tab:enroll} 显示，correct enrollment 与 shuffled、interferer、zero enrollment 之间存在巨大差距，说明 TD-SpeakerBeam 的输出由参考语音身份显著调制。这一结果对应认知意义上的选择性注意：系统根据“要听谁”的线索选择输出声源。

\subsubsection{从特征绑定到条件融合}

听觉对象形成要求系统把分散在时间和频率上的线索组织为同一说话者。SpeakerBeam 的适配机制可以理解为一种目标身份绑定：参考语音分支给出“目标是谁”，分离主干从 mixture 中寻找与该身份一致的语音成分。target-as-enrollment 略高于 correct enrollment，短 enrollment 和 noisy enrollment 逐步退化，也与这一解释一致：参考语音越接近目标身份且质量越高，绑定越稳定。表 \ref{tab:multi-enroll} 进一步显示，两条参考语音可以比单条参考提供更稳定的说话人线索，这对应人类听觉中由多次听闻形成的熟悉声纹记忆。

\subsubsection{从注意资源分配到高效推理}

认知系统的注意资源是有限的，计算模型也受到参数量、通道数和深度约束。表 \ref{tab:capacity} 显示，small student 到 mid student 的容量提升带来 2.27 dB \sisdr{} 增益；strong teacher 进一步提高到 12.52 dB。表 \ref{tab:stress} 进一步说明，当背景噪声增强或说话者数量从 2 人增加到 3 人时，固定容量模型的干扰抑制能力会明显下降。该现象说明目标说话人提取同时依赖参考语音线索和足够容量的分离主干，以建模长时上下文、目标音色和多干扰源抑制。蒸馏实验进一步表明，模型压缩需要与真实目标监督保持一致，教师输出更适合作为 fine-tune 正则。

\subsubsection{与空间声源定位的差异}

评分标准中的声源定位更接近多麦克风或双耳听觉建模问题，需要估计声源方向、距离或空间角度，并利用 ITD、ILD、HRTF、DOA 或麦克风阵列特征。本文系统没有实现这一部分。本文完成的是目标说话人条件分离：用 enrollment 表示“要听谁”，再从单通道混合语音中提取与该声纹匹配的声音。它与人类鸡尾酒会能力在选择性注意和特征绑定层面具有功能相似性，但缺少空间注意转移、左右耳线索整合和位置跟踪。后续若要覆盖声源定位，可将数据扩展为双通道或阵列录音，在模型输入中加入 IPD/ILD、角度标签或可学习 spatial embedding，并把目标身份条件与空间方向条件联合建模。

\subsection{附加挑战实验：流式处理、真实噪声与跨语言泛化}

本节作为任务二的补充实验，进一步考察系统在交互式展示、部署效率、真实环境噪声和跨语言场景下的表现。实验不改变主模型和主测试结论，只用于分析模型在更复杂条件下的适用边界。本节包含四个部分：交互式 Web 展示页面、流式实时处理模拟、DEMAND 真实环境噪声测试，以及中英跨语言泛化测试。

\subsubsection{交互式 Web 展示页面}

为直观展示目标说话人提取系统的工作过程，项目实现了一个基于 HTML/CSS/JavaScript 的交互式 Web 展示页面（\texttt{gui/index.html}）。展示页面包含四个视图：

\begin{itemize}
    \item \textbf{系统总览}：以信息图形式总结系统架构、核心指标（最终 SI-SDR 10.80 dB，单参考学生 10.59 dB）和实验设计理念。
    \item \textbf{案例展厅}：提供 3 组可交互的 demo 音频案例，每组包含原始 mixture、干净 target、基线输出和本文模型输出，用户可点击播放对比。每个案例附有说话人来源、难度标签、SI-SDR 和 SIR 等指标。
    \item \textbf{实验图表}：以雷达图和柱状图展示不同实验条件下的指标对比，包括主模型对比、容量对照、enrollment sanity check 和多参考聚合结果。
    \item \textbf{认知映射}：将鸡尾酒会问题的认知机制（选择性注意、特征绑定、资源分配）映射到 TD-SpeakerBeam 的计算组件（enrollment、speaker embedding、动态 block 路由）。
\end{itemize}

所有 demo 音频（\texttt{demo\_audio/} 目录）均来自 shared-clean80 测试集的实际输出，展示页面可离线运行。该页面服务于课程报告的辅助展示，不涉及后端或在线推理。

\subsubsection{流式实时处理模拟}

\paragraph{设计。}
流式处理脚本（\texttt{scripts/streaming\_demo.py}）以 chunk 为单位从 wav 文件读取混合语音，模拟实时多说话者语音分离场景。主要参数：chunk 2.0 秒（16000 采样点，$f_s = 8$ kHz），hop 0.5 秒（75\% 重叠），Hann 窗 overlap-add 拼接，参考语音全量预加载并复用每个 chunk，RMS 归一化与训练时 TSEDataset 保持一致。每个 chunk 依次经过归一化、模型推理和 overlap-add 重建，输出完整分离音频，同时记录每个 chunk 的延迟和实时率（RTF）。

\paragraph{结果。}
实验在 NVIDIA RTX 4080 SUPER 上运行。以 32 秒混合语音（8 段 4 秒测试样本拼接）为输入，共计 61 个 chunk。首个 chunk 包含 CUDA 内核初始化开销，延迟 0.235 s（RTF 0.117）；后续 60 个 chunk 延迟稳定在 5.3--7.2 ms（RTF 0.0026--0.0036）。表 \ref{tab:streaming} 汇总了整体流式性能。

\begin{table}[H]
\centering
\begin{tabular}{lrr}
\toprule
指标 & 冷启动（chunk 0） & 稳态（chunk 1--60） \\
\midrule
Chunk 延迟 & 0.235 s & 5.5 ms（均值） \\
Chunk RTF & 0.117 & 0.0028（均值） \\
\bottomrule
\midrule
\multicolumn{3}{l}{整体统计（61 chunks）} \\
\midrule
输入音频时长 & \multicolumn{2}{r}{32.0 s} \\
Chunk 数量 & \multicolumn{2}{r}{61} \\
总处理时间 & \multicolumn{2}{r}{0.577 s} \\
总 RTF & \multicolumn{2}{r}{0.018} \\
\bottomrule
\end{tabular}
\caption{流式处理性能（RTX 4080 SUPER，32 s 音频，61 chunks）。RTF $<$ 1 表示快于实时。}
\label{tab:streaming}
\end{table}

表 \ref{tab:streaming} 显示，模型在冷启动后处理速度约 350 倍于实时（稳态 RTF 0.0028），整体 RTF 0.018 即约 55 倍于实时，说明模型推理具备实时部署潜力。首个 chunk 延迟较高（0.235 s）主要来自 CUDA kernel 首次调用时的 JIT 编译开销，实际在线部署中可通过模型预热消除。需要指出的是，本实验为 wav 文件模拟流式输入，不包含真实麦克风采集、端点检测和持续会话管理，因此不能等同于完整的在线流式系统。

\subsubsection{DEMAND 真实环境噪声测试}

\paragraph{设计。}
噪声测试脚本（\texttt{scripts/test\_real\_noise.py}）采用两阶段设计。第一阶段对 =30$ 条测试样本进行 clean 推理，筛除 clean \sisdr{} $\leq 0$ dB 的异常样本（主要由 enrollment 过短或数据构造异常导致），保留 28 条。第二阶段对保留样本逐一施加 DEMAND 咖啡厅环境噪声，构造 3 种 SNR 条件：$+20$ dB、$+10$ dB、$+5$ dB。噪声数据来自 DEMAND（Diverse Environments Multichannel Acoustic Noise Database）PCAFETER 场景（CC BY 4.0），使用 ch01 声道，从 16 kHz 降采样至 8 kHz。对每种条件分别记录 \sisdr{}、\sisdr{}i、SDR、SIR 和 SAR，报告均值和中位数以检验极端值影响。全部数据来源、许可证和处理方式记录于 \texttt{data\_sources.md}。

\paragraph{结果。}

\begin{table}[H]
\centering
\small
\begin{tabular}{lrrrrrr}
\toprule
条件 & $n$ & \sisdr{} mean & \sisdr{} median & \sisdr{}i mean & SDR mean & SIR mean \\
\midrule
Clean & 28 & 11.86 & 12.04 & 12.21 & 11.85 & 23.32 \\
DEMAND $+20$ dB & 28 & 11.18 & 11.41 & 11.64 & 11.17 & 22.54 \\
DEMAND $+10$ dB & 28 & 7.64 & 8.04 & 8.92 & 7.65 & 16.86 \\
DEMAND $+5$ dB & 28 & 4.12 & 4.37 & 6.84 & 4.13 & 11.82 \\
\bottomrule
\end{tabular}
\normalsize
\caption{DEMAND 真实环境噪声鲁棒性测试（28 样本，筛选条件：clean \sisdr{} $>$ 0 dB）。报告均值和中位数；\sisdr{} 单位为 dB。}
\label{tab:noise}
\end{table}

表 \ref{tab:noise} 显示，加入 DEMAND 咖啡厅环境噪声后分离质量呈现清晰的单调下降。$+20$ dB 条件下仅下降 0.68 dB（均值），说明轻度环境噪声对模型影响很小；$+10$ dB 下降 4.22 dB；$+5$ dB 下降 7.74 dB，退化幅度随噪声增强而逐步增大。所有噪声条件下的 \sisdr{}i 均为正值（6.84--11.64 dB），确认模型在真实环境噪声下仍能从 mixture 中提取目标说话人并取得相对增益。中位数与均值非常接近（差异 $<$ 0.4 dB），说明结果未受极端值主导。

与主实验中 babble 噪声压力测试（表 \ref{tab:stress}）相比：$+10$ dB 下 babble 为 5.94 dB，DEMAND 为 7.64 dB；$+5$ dB 下 babble 为 3.36 dB，DEMAND 为 4.12 dB。DEMAND 条件下的绝对 \sisdr{} 略高于 babble，推测因为 babble 噪声包含更多人声干扰频率成分，而 DEMAND 咖啡厅噪声的频谱分布与目标语音重叠较少。需要指出，本实验筛选了 clean \sisdr{} $>$ 0 dB 的样本（28/30），筛选后的 clean 均值 11.86 dB 与主实验完整测试集 10.59 dB 基本可比。

\subsubsection{中英跨语言泛化测试}

\paragraph{设计。}
跨语言测试脚本（\texttt{scripts/test\_cross\_lingual.py}）支持两个方向的测试：（A）中文目标说话人 + 英文干扰说话人；（B）英文目标说话人 + 中文干扰说话人。每组需要 3 个 8 kHz 单声道 wav 文件（target、enrollment 来自同一说话人，interferer 来自不同说话人）。脚本构造 0 dB 等响度混合语音，运行推理并记录 \sisdr{} 和 \sisdr{}i。中文语音来自 THCHS-30（Apache 2.0），选取说话人 A23 和 B33；英文语音来自 LibriSpeech test-clean（CC BY 4.0），选取说话人 121 和 1089。所有音频统一转换为 8 kHz mono、峰值归一化至 0.95。

\paragraph{结果。}

\begin{table}[H]
\centering
\begin{tabular}{lrrrr}
\toprule
测试方向 & 目标语言 & 干扰语言 & \sisdr{}$\uparrow$ & \sisdr{}i$\uparrow$ \\
\midrule
中文目标 + 英文干扰 & zh & en & $-$33.03 & $-$32.99 \\
英文目标 + 中文干扰 & en & zh & $-$41.62 & $-$41.62 \\
\bottomrule
\end{tabular}
\caption{中英跨语言目标说话人提取结果。}
\label{tab:crosslingual}
\end{table}

表 \ref{tab:crosslingual} 显示两个方向的跨语言测试结果均为强负值，说明在本小规模中英跨语言压力测试中，模型出现完全失败。具体分析：

\begin{itemize}
    \item \textbf{中文目标方向}（\sisdr{} $-$33.03 dB）：enrollment 和 target 均为中文（同一说话人 A23），但模型仅在英文上训练，其 speaker embedding 分支从未学习过中文语音的声学特征，无法从 enrollment 中提取有效的目标声纹特征来指导分离。
    \item \textbf{英文目标方向}（\sisdr{} $-$41.62 dB）：enrollment 和 target 均为英文（同一说话人 121），模型具备英文声纹建模能力，但干扰说话人为中文语音——该干扰信号与模型见过的英文干扰声学分布不同，导致 masker 模块无法有效抑制。
\end{itemize}

两个方向的极端负值表明模型输出与目标几乎完全错位。这一结果说明，基于单语言（英文 LibriSpeech）训练的目标说话人提取模型\textbf{在跨语言场景下缺乏泛化能力}。要支持中英跨语言场景，需要在训练数据中引入多语言语音，使 speaker embedding 分支和 masker 分支同时学习跨语言的声纹表示和干扰建模。需要强调的是，本实验样本量极小（每个方向仅一对说话人），上述结论需要在更大规模的多语言测试中进一步验证。

\subsubsection{附加实验小结}

这三组附加实验和交互式 Web 展示共同说明，当前 TD-SpeakerBeam 系统在推理速度上具备实时部署潜力（32 s 音频 61 chunks，稳态 RTF 0.0028，整体 RTF 0.018），在 DEMAND 真实环境噪声下呈现清晰的 SNR 依赖退化规律（$+20$ dB 仅下降 0.68 dB，$+10$ dB 下降 4.22 dB，$+5$ dB 下降 7.74 dB，\sisdr{}i 均保持正值），但在跨语言场景下存在明显局限。跨语言测试表明，单一英文训练数据难以支撑中文目标说话人提取，后续需要引入多语言训练数据和更稳健的说话人表征。所有新增脚本、实验输出和数据来源文档均未修改原有训练与评估代码。

\subsection{结论}

本文实现并评估了一个基于 TD-SpeakerBeam 的目标说话人提取系统，用 mixture 与 enrollment 模拟鸡尾酒会场景中的选择性听觉。实验在 shared-clean80 数据上得到 10 dB 以上的稳定分离结果，strong teacher 达到 12.52 dB test \sisdr{}，mid student 达到 10.43 dB，fine-tune 蒸馏学生达到 10.59 dB，双参考语音 embedding-level 聚合进一步达到 10.80 dB。不同混合 SNR 和说话者数量实验表明，20 dB babble 条件下系统仍保持 9.77 dB \sisdr{}，但 0 dB babble 和 3-speaker clean 条件分别降至 0.91 dB 和 2.67 dB，说明复杂声场仍需要噪声增强训练或显式多说话者建模。参考语音变体实验显示，正确参考语音、错误参考语音、干扰说话人参考语音和零参考语音之间存在清晰区分，验证了系统的目标说话人选择能力。附加挑战实验进一步表明，模型在推理速度上具备实时部署潜力（32 s 音频 61 chunks，稳态 RTF 0.0028），在 DEMAND 真实噪声下 \sisdr{} 随 SNR 单调退化但 \sisdr{}i 始终为正（$+20$ dB 仅下降 0.68 dB，$+5$ dB 仍保持 4.12 dB），但单语言训练模型在中英跨语言压力测试中完全失败（\sisdr{} $< -30$ dB），为后续多语言 TSE 系统设计提供了明确方向。

本系统的主要局限是未显式建模空间声源定位，尚未利用双耳线索、DOA 或麦克风阵列信息；本文重点模拟目标驱动的选择性注意和声纹绑定。围绕数据构造、蒸馏和频域前端的诊断实验进一步给出三点经验：target/enrollment 必须严格 disjoint 并保证说话人标签正确；有干净目标监督时，从零波形蒸馏并不稳定；频域前端即使在训练阶段集成仍需谨慎调参，不能仅在推理阶段硬加。本文完成了从认知任务到可运行语音系统的建模、实现、评估和反思，为后续加入更稳健的参考语音增强、训练时频域前端、复杂度感知推理、多语言训练或双耳/阵列空间线索建模提供了实验基础。

\clearpage

\section{任务三：神经网络与反向传播的人类学习机制映射}

\subsection{算法选择与问题定位}

本文选择人工神经网络及其核心训练算法反向传播作为分析对象。这个算法适合放在“认知与计算”的语境下讨论：它广泛用于视觉、语音和自然语言处理；结构上借用了生物神经系统的启发；训练过程又围绕误差反馈展开，便于和人类学习中的反馈、修正和表征变化建立联系。

人工神经网络通过多层可调参数将输入映射到输出。给定输入 $x$、目标标签 $y$ 和模型预测 $\hat{y}=f_\theta(x)$，训练目标通常是最小化损失函数：

$$
\min_\theta L(f_\theta(x),y)
$$

反向传播算法利用链式法则计算损失函数对每一层参数的梯度，并通过梯度下降更新参数：

$$
\theta \leftarrow \theta - \eta \nabla_\theta L
$$

Rumelhart、Hinton 和 Williams 在 1986 年的经典论文中推广了通过反向传播误差学习内部表示的方法，使多层神经网络能够学习复杂映射\cite{rumelhart1986backprop}。这与联结主义强调的“连接权重随经验改变”相呼应，也成为后续深度学习模型的算法基础。

\subsection{心理学理论基础：联结主义、反馈学习与建构主义}

神经网络与人类学习之间的关系不宜只从“神经元”这一表面相似性理解，而应放在心理学和认知科学的多个理论传统中考察。首先是联结主义。联结主义认为，认知不由少量显式规则直接控制，而由大量简单单元之间的连接模式产生。知识以分布式方式存在，学习则表现为连接权重随经验而改变。人工神经网络正是联结主义在计算层面的重要形式。

其次是反馈学习。行为主义心理学强调外部反馈对行为塑造的作用，Thorndike 的效果律认为，带来满意结果的反应更容易被保留，带来不满意结果的反应更容易被削弱\cite{thorndike1911animal}。虽然反向传播不是行为主义模型，但它同样把“结果误差”作为学习更新的驱动力。模型预测错误越大，参数更新越明显；预测越接近目标，更新越小。这与人类在练习、纠错和反馈中逐渐改善表现具有功能上的相似性。

第三是建构主义学习观。建构主义强调学习者会在已有认知结构基础上主动建构理解。Piaget 认为儿童通过同化和顺应调整认知结构；Vygotsky 则强调社会互动和最近发展区对学习的重要性\cite{piaget1952origins,vygotsky1978mind}。神经网络并不具有人类儿童那样的主动建构和社会互动，但其表示学习也可以被理解为一种“内部结构重组”：模型在训练过程中不断改变中间层表征，使输入数据在新的表示空间中更容易被区分和预测。

神经网络与人类学习的对应关系并不是单线条的。它接近联结主义的分布式表示、反馈学习的误差修正和建构主义的表征重组，同时也暴露出主动探索、社会互动和具身经验不足的问题。

\begin{figure}[htbp]
\centering
\begin{tikzpicture}[
  font=\small,
  box/.style={draw=tseBlueLine, fill=tseBlue, rounded corners=4pt, align=center, minimum width=3.0cm, minimum height=0.78cm},
  feed/.style={draw=tseAmberLine, fill=tseAmber, rounded corners=4pt, align=center, minimum width=3.0cm, minimum height=0.78cm},
  update/.style={draw=tseGreenLine, fill=tseGreen, rounded corners=4pt, align=center, minimum width=3.7cm, minimum height=0.78cm},
  arrow/.style={-{Stealth[length=2.2mm]}, thick}
]
\node[box] (input) at (-4.8,1.1) {经验输入\\样本、任务、反馈};
\node[box] (repr) at (0,1.1) {内部表示\\特征、概念、图式};
\node[box] (pred) at (4.8,1.1) {行为/预测\\分类、生成、决策};
\node[feed] (error) at (4.8,-1.25) {误差信号\\损失、反馈、预测误差};
\node[update] (update) at (0,-1.25) {学习更新\\权重调整、策略修正、结构重组};
\draw[arrow] (input) -- (repr);
\draw[arrow] (repr) -- (pred);
\draw[arrow] (pred) -- (error);
\draw[arrow] (error) -- (update);
\draw[arrow] (update) -- (repr);
\end{tikzpicture}
\caption{神经网络学习与人类反馈学习的共同抽象结构}
\label{fig:learning_loop}
\end{figure}

图\ref{fig:learning_loop}展示了神经网络和人类反馈学习之间的共同抽象结构：经验输入形成内部表示，内部表示支持预测或行动，外部反馈产生误差信号，误差信号进一步改变内部结构。差异在于，神经网络的误差信号通常由数学损失函数定义，而人类学习中的反馈可能来自成绩、身体感觉、情绪体验、他人评价和社会规范。

\begin{longtable}{p{3.0cm}p{4.2cm}p{4.2cm}p{2.8cm}}
\toprule
机器学习概念 & 算法含义 & 心理学对应概念 & 主要差异 \\
\midrule
\endfirsthead
\toprule
机器学习概念 & 算法含义 & 心理学对应概念 & 主要差异 \\
\midrule
\endhead
神经元与连接权重 & 单元激活和参数化连接 & 神经元活动、突触强度、联结主义表征 & 人工神经元极度简化 \\
前向传播 & 输入逐层变换为输出 & 感知加工、证据累积、判断形成 & 人类加工包含反馈、注意和身体状态 \\
损失函数 & 衡量预测与目标差异 & 错误反馈、目标差距、任务评价 & 人类目标常常多元且变化 \\
反向传播 & 根据梯度分配误差并更新参数 & 反馈学习、试误学习、预测误差修正 & 标准反传缺乏生物可行性 \\
表示学习 & 中间层形成抽象特征 & 概念形成、图式更新、类别学习 & 人类概念含情境和价值意义 \\
迁移学习 & 复用已有表示解决新任务 & 经验迁移、类比学习 & 人类能主动判断迁移条件 \\
课程学习 & 从简单样本过渡到复杂样本 & 由易到难、最近发展区 & 人类课程还依赖动机和社会支架 \\
知识蒸馏 & 小模型模仿强模型输出 & 师徒学习、示范学习、模仿学习 & 人类模仿包含意图理解 \\
\bottomrule
\caption{神经网络算法与人类学习心理过程的对应关系}
\label{tab:nn_psychology_mapping}\\
\end{longtable}

\subsection{神经网络与联结主义认知观}

联结主义认为，心智并不一定依赖显式符号规则，而可以由许多简单处理单元的并行活动产生。人工神经网络正是这种思想的计算实现：每个神经元执行加权求和和非线性变换，多层神经元组合形成复杂函数。知识并不存储在某个单独规则中，而分布在网络连接权重中。

这种分布式表示与人类认知有一定相似性。例如，人类识别“猫”不依赖单一特征，而会综合耳朵形状、眼睛、轮廓、运动方式、叫声和以往经验。神经网络也通过多层特征逐级组合，从边缘、纹理等低层特征，到部件、形状和语义概念等高层特征。这种层级表示使深度学习在图像识别和语音识别中表现突出。

但是，人工神经网络与真实大脑仍有显著差异。生物神经元具有复杂的时序放电、突触可塑性、神经递质调节和反馈环路，而人工神经元通常是高度简化的数学函数。人工网络的训练往往依赖大量标注数据和全局损失函数，而人类学习更多依赖多模态经验、动机、注意、社会反馈和身体行动。

\subsection{反向传播与人类学习中的反馈机制}

反向传播与人类学习最直接的相似之处是二者都依赖反馈。学生做题后通过答案知道自己错在哪里，运动员通过教练反馈调整动作，儿童通过成人纠正学习语言表达。这些过程都包含“预测—反馈—修正”的循环。反向传播也可以看作机器学习中的误差修正机制：模型先给出预测，再根据预测与真实标签之间的误差调整内部参数。

这种机制与心理学中的试误学习、强化学习和误差驱动学习具有共同结构。Thorndike 的效果律认为，带来满意结果的行为更可能被保留，带来不满意结果的行为更可能被削弱。行为主义强调外部反馈对行为塑造的作用。反向传播虽然不是行为主义模型，但也体现了结果误差对系统内部结构调整的影响。

在认知神经科学中，预测误差也是重要概念。大脑不断根据当前模型预测外部输入，当预测与实际输入不一致时，误差信号促使内部模型更新。这与深度学习中的损失函数和梯度更新存在类比关系。不同的是，人类预测误差通常是局部、多模态、连续且与注意和情绪相关的，而反向传播中的误差信号是数学上精确定义的全局目标函数梯度。

\subsection{神经网络与 Hebbian 学习的关系}

Hebbian 学习通常概括为“同时激活的神经元连接会增强”。其基本思想是，如果两个神经元经常共同活动，它们之间的连接权重会增加\cite{hebb1949organization}。这一规则强调局部相关性，不需要全局损失函数，因此更接近生物突触可塑性。

反向传播与 Hebbian 学习存在明显差异。Hebbian 学习是局部的、无监督或弱监督的，强调共同激活；反向传播是全局误差驱动的，需要从输出层向前层传播梯度。近年来，一些研究尝试将 Hebbian 规则与卷积神经网络结合，以提高模型的生物合理性。例如 2025 年的研究探索了 Hebbian 卷积网络、竞争机制和 BCM 学习规则，试图在图像任务中获得更接近生物可行性的表示学习能力\cite{nimmo2025hebbian}。

因此，反向传播可以被看作工程上高效的学习算法，但不完全是生物大脑的直接模拟。它与人类学习的关系更适合理解为\textbf{功能类比}而非\textbf{机制等同}：二者都利用反馈减少错误，但大脑不太可能以标准反向传播的形式精确计算每个突触的全局梯度。

\subsection{机器学习机制与心理学习机制的主要对照}

神经网络不是单一心理学理论的机械翻版，它把多种学习机制抽象为可计算形式。表\ref{tab:learning_mechanism_summary}按“机器学习机制—心理学来源—解释力—局限性”归纳这些对应关系。可以看到，机器学习借用了人类学习中的一些可形式化片段，例如误差反馈、经验迁移、由易到难训练和示范学习；这些机制进入工程系统后，会被转化为优化、表示和控制问题。

\begin{longtable}{p{3.0cm}p{3.4cm}p{4.2cm}p{3.4cm}}
\toprule
机器学习机制 & 心理学/认知科学来源 & 能解释的人类学习侧面 & 主要局限 \\
\midrule
\endfirsthead
\toprule
机器学习机制 & 心理学/认知科学来源 & 能解释的人类学习侧面 & 主要局限 \\
\midrule
\endhead
神经网络表示学习 & 联结主义、分布式表征 & 概念由多维特征组合形成，知识不一定是显式规则 & 缺少具身和社会语义 \\
反向传播 & 误差驱动学习、预测误差 & 从错误中修正内部表示，提高任务表现 & 需要全局梯度，生物可行性有限 \\
Hebbian/局部学习 & 突触可塑性、关联学习 & 共同出现的经验会增强心理联结 & 难以单独解决复杂监督任务 \\
强化学习 & 效果律、奖惩学习 & 行为结果影响后续选择 & 人类奖惩包含价值和社会意义 \\
迁移学习 & 类比、经验复用 & 旧知识帮助新任务学习 & 机器迁移容易受分布相似性限制 \\
课程学习 & 由易到难、最近发展区 & 学习顺序影响掌握效率和动机 & 机器课程通常缺少情绪和社会支架 \\
知识蒸馏 & 模仿学习、师徒学习 & 弱学习者向强示范者学习 & 蒸馏多传递输出，不传递意图和解释 \\
持续学习 & 终身学习、记忆巩固 & 新旧经验不断整合 & 神经网络仍易灾难性遗忘 \\
\bottomrule
\caption{机器学习机制与心理学习机制的综述式归纳}
\label{tab:learning_mechanism_summary}\\
\end{longtable}

\subsection{神经网络如何模拟人类学习过程}

\subsubsection{从感知特征到抽象概念}

深度神经网络的层级结构可以模拟人类从低层感知到高层概念的加工过程。在视觉任务中，浅层卷积核通常学习边缘、方向和颜色等局部特征，中层学习纹理和部件，高层学习物体类别。这与人类视觉系统中从初级视觉皮层到高级视觉区域的层级加工有一定相似性。

在语言模型中，浅层可能更多处理词形和局部句法，高层则捕捉语义关系、篇章结构和任务意图。虽然这种对应不是严格的一一映射，但说明神经网络可以通过多层表示逐步形成抽象能力。

\subsubsection{通过反馈调整内部表示}

神经网络在训练中不断根据误差调整权重，使内部表示更适合任务目标。人类学习也会根据反馈调整认知策略。例如，学生一开始可能只记住公式，后来通过错题反馈逐渐理解适用条件和解题结构。二者都表现为内部表征随经验更新。

\subsubsection{迁移学习与经验复用}

预训练神经网络可以将一个任务中学到的表示迁移到另一个任务。例如，图像模型在大数据集上学到边缘、纹理和形状表示后，可以迁移到医学影像或工业缺陷检测任务。人类也具有迁移学习能力，例如学过羽毛球的人学习网球时会迁移身体协调、击球时机和空间判断经验。

人类迁移更灵活，能够结合类比、解释和常识主动选择迁移哪些知识；神经网络的迁移则依赖表示空间相似性和微调策略，可能出现负迁移。

\subsubsection{课程学习与由易到难}

课程学习是神经网络训练中与人类学习高度相似的策略。人类学习通常从简单任务开始，再逐步进入复杂任务。例如学习数学时先学加减乘除，再学方程、函数和微积分；学习语言时先掌握词汇和短句，再学习复杂语法和修辞。机器学习中的 curriculum learning 同样通过控制样本顺序，使模型先学习低难度样本，再逐步接触高难度样本。Bengio 等人在 2009 年系统提出课程学习思想，指出合理的样本排序可能改善模型优化和泛化\cite{bengio2009curriculum}。

结合本项目可以看到，课程学习思想很适合解释语音分离训练的难度设计：模型可以先学习低重叠、低噪声、两人语音，再逐步接触高重叠、强噪声、三人语音和短参考语音条件。这种设计把“由易到难”的学习经验转化成了样本调度策略。

\subsection{神经网络如何模拟人类决策过程}

神经网络的决策过程通常表现为输入特征经过多层变换后得到输出概率。例如分类模型会输出每个类别的概率分布，选择概率最高的类别作为决策。这与人类决策中的证据累积过程有一定相似性。人类在识别物体或判断语义时，也会综合多个线索并形成置信度。

在心理学中，许多决策模型认为人类会在不确定条件下积累证据，当某一选择的证据超过阈值时作出判断。神经网络中的 softmax 输出虽然不是严格的心理过程模型，但可以视作一种置信度表达。模型训练中通过交叉熵损失使正确类别概率提高，类似于经验不断强化正确判断路径。

人类决策还会受到动机、情绪、价值、社会规范和长期目标影响。一个人可能明知某选择短期收益较低，但基于道德或长期规划仍选择它。神经网络通常只优化给定目标函数，缺乏自主设定目标和价值权衡能力。因此，神经网络能够模拟\textbf{感知分类和模式判断}，但难以完整模拟\textbf{价值驱动的人类决策}。

\subsection{案例分析：从儿童概念学习到语音分离模型训练}

为了更具体地理解神经网络与人类学习的映射，可以比较儿童学习动物概念和语音分离模型学习目标说话人的过程。儿童学习“猫”这一概念时，最初可能依靠毛茸茸、会叫、会动等感性特征；随着经验增加，儿童逐渐能区分猫、狗、玩具猫和卡通猫，并能在不同姿态和环境中识别猫。这一过程包含多模态经验、社会反馈、概念修正和语言标签。

神经网络学习图像分类或语音分离时，也经历从低层特征到高层表征的过程。以语音分离为例，模型首先需要处理波形或频谱中的局部模式，例如音高、共振峰、能量变化和时间包络；随后学习说话人身份、语音片段连续性和目标声源与干扰声源之间的差异。在目标说话人提取任务中，enrollment 语音相当于给模型一个“要听谁”的条件线索，模型根据这个线索从混合语音中提取目标说话人。这与人类在鸡尾酒会环境中根据熟悉声音、语义线索或注意目标持续跟踪某个人的声音具有功能相似性。

二者差异也很明显。儿童并不需要听到上万次标注样本才能认识一个人，也不会只根据某个固定损失函数学习。人类可以主动询问、观察表情、利用语境推断身份，并在错误后通过社会反馈修正理解。语音分离模型则依赖大量 mixture-target 训练对，其学习目标由研究者定义，泛化能力受数据分布、模型结构和训练策略限制。本项目中，自写模型失败而开源 SpeakerBeam 成功的过程也说明：神经网络学习不是有数据就自然成功，合适的结构、可靠的条件建模和严格的数据审计同样重要。

知识蒸馏可以从心理学角度理解为一种“向教师学习”的机制。强教师模型先学习到较好的输入输出映射，小型学生模型再同时学习真实标签和教师输出，以较小参数量逼近教师行为。它类似学生通过模仿教师解题思路来提高效率，但机器蒸馏主要传递输出分布或波形目标，不传递教师的意图、解释、情绪鼓励和学习策略。在本文实验中，从零蒸馏效果低于监督训练，也提醒我们不能把“向教师学习”简单套用为必然有效的工程结论。

\subsection{神经网络未能模拟的人类学习特征}

神经网络与人类学习仍有几个明显差距。人类可以通过少量示例快速形成概念，而神经网络通常需要大量数据；大模型的少样本能力也主要来自大规模预训练。人类学习还带有主动性，儿童会探索环境、提出问题、选择关注对象，神经网络通常只是按照外部给定的数据和目标训练。人类概念建立在身体经验和社会互动中，理解“重”“痛”“冷”“危险”等概念涉及感官和情绪体验；神经网络更多学习词语和特征之间的统计关联。人类还能长期学习新知识并保持旧知识，神经网络在连续学习中仍容易出现灾难性遗忘。近年来也有研究探索不依赖反向传播的 Hebbian 反馈网络，用局部规则模拟持续学习中的表征保留\cite{li2026hebbiancontinual}，说明机器学习仍在寻找更接近生物学习的机制。

\subsection{批判性讨论：相似不等于等同}

神经网络与人类学习之间存在许多有价值的类比：神经元与神经元、权重与突触强度、误差反馈与学习反馈、层级表示与感知层级、课程学习与由易到难训练。但这些类比不能被过度解释。

首先，人工神经元极度简化，无法反映真实神经元的复杂生物物理过程。其次，反向传播需要全局误差信号和精确权重更新，而大脑更可能依赖局部可塑性、反馈调制和多种神经递质机制。再次，神经网络学习通常目标单一，而人类学习包含动机、情绪、社会评价和意义建构。

更重要的是，机器学习中的“学习”往往是任务性能优化，而心理学中的“学习”还包括主体经验的改变。一个学生学习一门课程，不仅获得解题能力，也可能改变自我效能感、兴趣、职业想象和社会关系。一个儿童学习语言，不只是拟合词语分布，而是在共同注意、模仿、游戏、情绪回应和文化规范中成为社会成员。相比之下，神经网络参数更新虽然可以提高任务表现，却没有自我经验，也不会理解学习对自身生命历程的意义。

神经网络适合作为人类学习的局部计算模型，而不是完整心理模型。它可以帮助解释表征如何通过经验改变、反馈如何改善预测、复杂模式如何由简单单元组合产生；若要解释人类学习的全部面貌，还需要结合发展心理学、教育心理学、社会文化理论、强化学习、具身认知和神经可塑性研究。

从认知与计算课程的角度看，这种“不完全对应”本身具有启发意义。计算模型的价值不在于完全复制人类，而在于把某些认知机制抽象成可实现、可测试的形式。反向传播把“从错误中学习”变成数学优化，课程学习把“由易到难”变成样本调度，知识蒸馏把“向教师学习”变成模型压缩。这些抽象会丢失许多人类学习的丰富性，但也使认知机制能够进入工程系统并接受实验检验。

\subsection{小结}

神经网络及反向传播算法在一定程度上模拟了人类学习中的反馈修正、层级表征、经验迁移和由易到难学习过程。它们体现了联结主义认知观，即复杂智能可以由大量简单单元的连接权重变化涌现出来。然而，标准神经网络仍难以模拟人类学习中的小样本概念形成、主动探索、具身语义、情绪动机和终身学习。反向传播与人类学习之间更应被理解为功能上的相似，而非生物机制上的等同。

本节的核心结论是：神经网络可以把人类学习中的部分机制转化为可计算形式，例如把反馈形式化为误差、把经验形式化为样本、把概念形式化为表示、把教师示范形式化为蒸馏目标。但人类学习还包含主体经验、社会意义、身体基础和价值方向。未来更具认知合理性的人工智能，可能需要把深度学习、Hebbian 可塑性、预测加工、强化学习、具身交互和社会学习机制结合起来。

\addcontentsline{toc}{section}{参考文献}
\begingroup
\zihao{-5}\songti
\begin{thebibliography}{99}
\bibitem{turing1950computing} Turing, A. M. Computing machinery and intelligence. \textit{Mind}, 1950, 59(236): 433--460.
\bibitem{newell1972human} Newell, A., Simon, H. A. \textit{Human Problem Solving}. Prentice-Hall, 1972.
\bibitem{varela1991embodied} Varela, F. J., Thompson, E., Rosch, E. \textit{The Embodied Mind: Cognitive Science and Human Experience}. MIT Press, 1991.
\bibitem{clark1997being} Clark, A. \textit{Being There: Putting Brain, Body, and World Together Again}. MIT Press, 1997.
\bibitem{openai2023gpt4} OpenAI. GPT-4 technical report. \textit{arXiv preprint} arXiv:2303.08774, 2023.
\bibitem{alphafold2024nature} Abramson, J., Adler, J., Dunger, J., et al. Accurate structure prediction of biomolecular interactions with AlphaFold 3. \textit{Nature}, 2024, 630(8016): 493--500.
\bibitem{nori2023gpt4medical} Nori, H., King, N., McKinney, S. M., Carignan, D., Horvitz, E. Capabilities of GPT-4 on medical challenge problems. \textit{arXiv preprint} arXiv:2303.13375, 2023.
\bibitem{ying2025benchmarking} Ying, L., Collins, K. M., Wong, L., et al. On benchmarking human-like intelligence in machines. \textit{arXiv preprint} arXiv:2502.20502, 2025.
\bibitem{cherry1953} Cherry, E. C. Some experiments on the recognition of speech, with one and with two ears. \textit{The Journal of the Acoustical Society of America}, 1953, 25(5): 975--979.
\bibitem{broadbent1958} Broadbent, D. E. \textit{Perception and Communication}. Pergamon Press, 1958.
\bibitem{bregman1990} Bregman, A. S. \textit{Auditory Scene Analysis: The Perceptual Organization of Sound}. MIT Press, 1990.
\bibitem{hershey2016deep} Hershey, J. R., Chen, Z., Le Roux, J., Watanabe, S. Deep clustering: Discriminative embeddings for segmentation and separation. \textit{Proceedings of IEEE ICASSP}, 2016: 31--35.
\bibitem{yu2017pit} Yu, D., Kolb{\ae}k, M., Tan, Z.-H., Jensen, J. Permutation invariant training of deep models for speaker-independent multi-talker speech separation. \textit{IEEE/ACM Transactions on Audio, Speech, and Language Processing}, 2017, 25(10): 2066--2079.
\bibitem{luo2019convtasnet} Luo, Y., Mesgarani, N. Conv-TasNet: Surpassing ideal time-frequency magnitude masking for speech separation. \textit{IEEE/ACM Transactions on Audio, Speech, and Language Processing}, 2019, 27(8): 1256--1266.
\bibitem{luo2020dprnn} Luo, Y., Chen, Z., Yoshioka, T. Dual-path RNN: Efficient long sequence modeling for time-domain single-channel speech separation. \textit{Proceedings of IEEE ICASSP}, 2020: 46--50.
\bibitem{subakan2021sepformer} Subakan, C., Ravanelli, M., Cornell, S., Bronzi, M., Zhong, J. Attention is all you need in speech separation. \textit{Proceedings of IEEE ICASSP}, 2021: 21--25.
\bibitem{wang2023tfgridnet} Wang, Z.-Q., Cornell, S., Choi, S., Lee, Y., Kim, B.-Y., Watanabe, S. TF-GridNet: Integrating full- and sub-band modeling for speech separation. \textit{IEEE/ACM Transactions on Audio, Speech, and Language Processing}, 2023, 31: 3221--3236.
\bibitem{wang2019voicefilter} Wang, Q., Muckenhirn, H., Wilson, K., et al. VoiceFilter: Targeted voice separation by speaker-conditioned spectrogram masking. \textit{Proceedings of Interspeech}, 2019.
\bibitem{zmolikova2019speakerbeam} Zmol{\'i}kov{\'a}, K., Delcroix, M., Kinoshita, K., Ochiai, T., Nakatani, T., Burget, L., \v{C}ernock{\'y}, J. SpeakerBeam: Speaker aware neural network for target speaker extraction in speech mixtures. \textit{IEEE Journal of Selected Topics in Signal Processing}, 2019, 13(4): 800--814.
\bibitem{ge2020spexplus} Ge, M., Xu, C., Wang, L., Chng, E. S., Li, H. SpEx+: A complete time domain speaker extraction network. \textit{Proceedings of Interspeech}, 2020.
\bibitem{hinton2015distill} Hinton, G., Vinyals, O., Dean, J. Distilling the knowledge in a neural network. \textit{arXiv preprint} arXiv:1503.02531, 2015.
\bibitem{vincent2006bss} Vincent, E., Gribonval, R., F{\'e}votte, C. Performance measurement in blind audio source separation. \textit{IEEE Transactions on Audio, Speech and Language Processing}, 2006, 14(4): 1462--1469.
\bibitem{leroux2019sdr} Le Roux, J., Wisdom, S., Erdogan, H., Hershey, J. R. SDR--half-baked or well done? \textit{Proceedings of IEEE ICASSP}, 2019: 626--630.
\bibitem{rumelhart1986backprop} Rumelhart, D. E., Hinton, G. E., Williams, R. J. Learning representations by back-propagating errors. \textit{Nature}, 1986, 323: 533--536.
\bibitem{thorndike1911animal} Thorndike, E. L. \textit{Animal Intelligence: Experimental Studies}. Macmillan, 1911.
\bibitem{hebb1949organization} Hebb, D. O. \textit{The Organization of Behavior: A Neuropsychological Theory}. Wiley, 1949.
\bibitem{piaget1952origins} Piaget, J. \textit{The Origins of Intelligence in Children}. International Universities Press, 1952.
\bibitem{vygotsky1978mind} Vygotsky, L. S. \textit{Mind in Society: The Development of Higher Psychological Processes}. Harvard University Press, 1978.
\bibitem{bengio2009curriculum} Bengio, Y., Louradour, J., Collobert, R., Weston, J. Curriculum learning. \textit{Proceedings of ICML}, 2009: 41--48.
\bibitem{nimmo2025hebbian} Nimmo, J. J., Mondragon, E. Advancing the biological plausibility and efficacy of Hebbian convolutional neural networks. \textit{arXiv preprint} arXiv:2501.17266, 2025.
\bibitem{li2026hebbiancontinual} Li, J. A backpropagation-free feedback-Hebbian network for continual learning dynamics. \textit{arXiv preprint} arXiv:2601.06758, 2026.
\bibitem{thiemann2013demand} Thiemann, J., Ito, N., Vincent, E. DEMAND: Diverse Environments Multichannel Acoustic Noise Database. \textit{Zenodo}, 2013. DOI: 10.5281/zenodo.1227121.
\bibitem{wang2015thchs30} Wang, D., Zhang, X. THCHS-30: A Free Chinese Speech Corpus. \textit{arXiv preprint} arXiv:1512.01882, 2015.
\end{thebibliography}
\endgroup

\end{document}
